\def\draft{n}
\documentclass[12pt]{amsart}
\pdfoutput=1
%\usepackage{showkeys}
\usepackage[headings]{fullpage}
\usepackage{amssymb,epic,eepic,epsfig}%,amsbsy,amsmath,bbold}
%\usepackage{pb-diagram,lamsarrow,pb-lams,amscd}
\usepackage{graphicx}
\usepackage{texdraw}
\usepackage{url}
\usepackage{subfigure}
\usepackage{multirow}
\usepackage{array}
%\usepackage{palatino} % use also: \documentclass[10.5pt]{amsart}
\usepackage[bookmarks=true,%
    colorlinks=true,%
    linkcolor=blue,%
    citecolor=blue,%
    filecolor=blue,%
    menucolor=blue,%
    urlcolor=blue,%
    breaklinks=true]{hyperref}

\usepackage{enumerate}
\usepackage{enumitem}
  
\usepackage{scalerel} 

\usepackage{mathdots}
\usepackage{amsfonts}

\usepackage{caption}
\usepackage{capt-of}

\usepackage[all]{xy}


%%%%%%%%%%% macros

\newtheorem{theorem}{Theorem}[section]
\theoremstyle{definition}
\newtheorem{proposition}[theorem]{Proposition}
\newtheorem{lemma}[theorem]{Lemma}
\newtheorem{definition}[theorem]{Definition}
\newtheorem{convention}[theorem]{Convention}
\newtheorem{remark}[theorem]{Remark}
\newtheorem{corollary}[theorem]{Corollary}
\newtheorem{conjecture}[theorem]{Conjecture}
\newtheorem{question}[theorem]{Question}
\newtheorem{problem}[theorem]{Problem}
\newtheorem{example}[theorem]{Example}
\newtheorem{exercise}[theorem]{Exercise}
\newtheorem{hint}[theorem]{Hint}
\newtheorem{notation}[theorem]{Notation}
\newtheorem{observation}[theorem]{Observation}


\numberwithin{equation}{section}

\def\printname#1{
        \if\draft y
                \smash{\makebox[0pt]{\hspace{-0.5in}
                        \raisebox{8pt}{\tt\tiny #1}}}
        \fi
}

\newcommand{\mathmode}[1]{$#1$}
\newcommand{\psdraw}[2]
         {\begin{array}{c} \hspace{-1.3mm}
        \raisebox{-4pt}{\epsfig{figure=draeq:FigureSelfEdgews/#1.eps,width=#2}}
        \hspace{-1.9mm}\end{array}}
\def\ppsdraw#1#2{\epsfig{figure=draws/#1.eps,height=#2in,silent=}}

\newlength{\standardunitlength}
\setlength{\standardunitlength}{0.0125in}
\newcommand{\eepic}[2]
        {\printname{#1}
        \setlength{\unitlength}{#2\standardunitlength}
        \begin{array}{c}  \hspace{-1.7mm}
                \raisebox{-8pt}{\input draws/#1.eepic }
                \hspace{-1.9mm}
        \end{array}
}
\newcommand{\silenteepic}[2]
        {\setlength{\unitlength}{#2\standardunitlength}
        \begin{array}{c}  \hspace{-1.7mm}
                \raisebox{-3pt}{\input draws/#1.tex }
                \hspace{-1.9mm}
        \end{array}
}

\catcode`\@=11
\long\def\@makecaption#1#2{%
     \vskip 10pt

\setbox\@tempboxa\hbox{%\ifvoid\tinybox\else\box\tinybox\fi
       \small\sf{\bfcaptionfont #1. }\ignorespaces #2}%
     \ifdim \wd\@tempboxa >\captionwidth {%
         \rightskip=\@captionmargin\leftskip=\@captionmargin
         \unhbox\@tempboxa\par}%
       \else
         \hbox to\hsize{\hfil\box\@tempboxa\hfil}%
     \fi}
\font\bfcaptionfont=cmssbx10 scaled \magstephalf
\newdimen\@captionmargin\@captionmargin=2\parindent
\newdimen\captionwidth\captionwidth=\hsize
\catcode`\@=12

\newcommand{\End}{\operatorname{End}}
\newcommand{\Span}{\operatorname{Span}}
\newcommand{\tr}{\operatorname{tr}}
\newcommand{\im}{\operatorname{Im}}

\newcommand{\Wtilde}[1]{\stackrel{\sim}{\smash{#1}\rule{0pt}{1.1ex}}}
\newcommand{\What}[1]{\stackrel{\widehat{}}{\smash{#1}\rule{0pt}{1.1ex}}}


\def\lbl#1{\label{#1}\printname{#1}}
\def\biblbl#1{\bibitem{#1}\printname{#1}}

\def\eqdef{\overset{\text{def}}{=}}

\def\BTclear{\def\tensorname{}
        \def\basea{}\def\collectora{}\def\emittera{}
        \def\baseb{}\def\collectorb{}\def\emitterb{}}

\newcommand{\prtag}[1]{\tag*{\llap{$#1$\hskip-\displaywidth}}}

\def\truncsimp{\overline{\Delta}}
\def\dtruncsimp{\overline{\overline{\Delta}}}
\def\cxymatrix#1{\xy*[c]\xybox{\xymatrix#1}\endxy}

\def\Mtildehat{\scalerel*{\widehat{\widetilde M}}{\widetilde{M^2}}}

%%%%%%%%%%%%% Stavros's defs

\def\BN{\mathbb N}
\def\BZ{\mathbb Z}
\def\BP{\mathbb P}
\def\BQ{\mathbb Q}
\def\BR{\mathbb R}
\def\BC{\mathbb C}
\def\BE{\mathbb E}
\def\BT{\mathbb T}
\def\BK{\mathbb K}

\def\calA{\mathcal A}
\def\calB{\mathcal B}
\def\calC{\mathcal C}
\def\calD{\mathcal D}
\def\D{\Delta}
\def\calK{\mathcal K}
\def\calL{\mathcal L}
\def\calG{\mathcal G}
\def\calO{\mathcal O}
\def\calT{\mathcal T}
\def\Pt{P}

\def\a{\alpha}
\def\bb{\beta}
\def\La{\Lambda}
\def\l{\lambda}
\def\Ga{\Gamma}
\def\S{\Sigma}
\def\s{\sigma}
\def\varsi{\varsigma}
\def\ga{\gamma}

\def\la{\langle}
\def\ra{\rangle}


\def\w{\omega}


\def\we{\wedge}

\def\g{\gamma}
\def\t{\tau}
\def\e{\epsilon}
\def\Ga{\Gamma}
\def\d{\delta}
\def\b{\beta}
\def\th{\theta}
\def\Th{\Theta}
\def\s{\sigma}
\def\p{\prime}

\def\rk{\mathrm{rk}}

\def\ti{\widetilde}

\def\calL{\mathcal L}
\def\longto{\longrightarrow}
\def\w{\omega}
\def\SL{\mathrm{SL}}
\def\PSL{\mathrm{PSL}}
\def\PGL{\mathrm{PGL}}
\def\GL{\mathrm{GL}}
\def\pt{\partial}
\def\Om{\Omega}
\def\om{\omega}
\def\calB{\mathcal{B}}
\def\z{\zeta}
\def\om{\omega}
\def\Om{\Omega}
\def\ep{\epsilon}
\def\sgn{\mathrm{sgn}}
\def\calR{\mathcal R}

\DeclareMathOperator{\Conj}{Conj}
\DeclareMathOperator{\id}{id}
\DeclareMathOperator{\sub}{sub}

\def\bS{\bar S}
\def\bs{\hat s}
\def\sign{\operatorname{sign}}
\def\Zp{\BZ_{\succ}[q] }
\def\fsl{\mathfrak{sl}}
\def\wp{w_{\succ}}
\def\wlt{w_{\mathrm{lt}}}
\def\mindeg{\mathrm{mindeg}}
\def\coeff{\mathrm{coeff}}

\def\be{  \begin{equation} }
\def\ee{  \end{equation} }
\def\cP{\mathcal P}
\def\A{\mathcal A}
\def\ve{\varepsilon}
\def\U{\mathcal U}
\def\Hom{\mathrm{Hom}}

\def\BH{\mathbb H}
\def\diag{\text{diag}}
\def\hb{\hbar}
\def\Sp{\mathrm{Sp}}

\def\can{\mathrm{can}}
\def\vol{\mathrm{Vol}}
\def\geom{\mathrm{geom}}
\def\bs{$\clubsuit$}
\def\Li{\mathrm{Li}}
\def\CT{\mathcal T}


\def\R{\mathbb{R}}
\def\Z{\mathbb{Z}}
\def\C{\mathbb{C}}
\def\N{\mathbb{N}}
\def\Q{\mathbb{Q}}
\def\T{\mathcal{T}}

\def\Log{\mathrm{Log}}
\def\rd{\mathrm{red}}
\def\Tr{\mathrm{Tr}}
\def\Ker{\mathrm{Ker}}
\def\Imag{\mathrm{Im}}
\def\Coker{\mathrm{Coker}}
\def\rank{\mathrm{rank}}
\def\Span{\mathrm{Span}}
\def\Stab{\mathrm{Stab}}
\usepackage{ifpdf}
\ifpdf
\DeclareGraphicsRule{*}{mps}{*}{}
\fi

\begin{document}

\begin{abstract}
The Ptolemy coordinates for boundary-unipotent $\SL(n,\C)$-representations of a 
3-manifold group were introduced in 
\cite{GaroufalidisThurstonZickert} inspired by the $\mathcal A$-coordinates 
on higher Teichm\"{u}ller space due to Fock and Goncharov. In this paper, 
we define the Ptolemy field of a (generic) $\PSL(2,\C)$-representation and prove that it coincides with the trace field
of the representation. This gives an efficient algorithm to compute the
trace field of a cusped hyperbolic manifold.%, which we implement and illustrate
%with numerous examples.
\end{abstract}


%%%%%%%%%%%%%%%%%%%%%%{page1}

\title[The Ptolemy field of $3$-manifold-representations]{
The Ptolemy field of $3$-manifold-representations}
\author{Stavros Garoufalidis}
\address{School of Mathematics \\
         Georgia Institute of Technology \\
         Atlanta, GA 30332-0160, USA \newline
         {\tt \url{http://www.math.gatech.edu/~stavros }}}
\email{stavros@math.gatech.edu}
\author{Matthias Goerner}
\address{University of Maryland \\
         Department of Mathematics \\
         College Park, MD 20742-4015, USA \newline
         {\tt \url{http://math.berkeley.edu/~matthias/}}}
\email{enischte@gmail.com}
\author{Christian K. Zickert}
\address{University of Maryland \\
         Department of Mathematics \\
         College Park, MD 20742-4015, USA \newline
         {\tt \url{http://www2.math.umd.edu/~zickert}}}
\email{zickert@umd.edu}
\thanks{S.~G.~and C.~Z.~were supported in part by the NSF. \\
\newline
1991 {\em Mathematics Classification.} Primary 57N10. Secondary 57M27.
\newline
{\em Key words and phrases: Ptolemy coordinates, trace field, 3-manifolds.
}
}

\date{\today }

\maketitle

\tableofcontents

\section{Introduction}

\subsection{The Ptolemy coordinates}
\label{sub.ge}

The Ptolemy coordinates for boundary-unipotent representations of a 
3-manifold group in $\SL(n,\C)$ were introduced in 
Garoufalidis--Thurston--Zickert~\cite{GaroufalidisThurstonZickert} inspired by the $\mathcal A$-coordinates 
on higher Teichm\"{u}ller space due to Fock and Goncharov
\cite{FockGoncharov}. In this paper we will focus primarily on representations in $\SL(2,\C)$ and $\PSL(2,\C)$.  %We begin 
%with an elementary introduction to this subject.

Given a topological ideal triangulation $\T$ of an oriented compact 3-manifold $M$, 
a \emph{Ptolemy assignment} (for $\SL(2,\C)$) is an assignment of a non-zero 
complex number (called a \emph{Ptolemy coordinate}) to each $1$-cell of $\CT$, such that for each simplex, the Ptolemy coordinates assigned to the edges $\varepsilon_{ij}$
satisfy the \emph{Ptolemy relation}
\begin{equation}\label{eq:PtolemyRelation}
c_{03}c_{12}+c_{01}c_{23}=c_{02}c_{13}.
\end{equation}
The set of Ptolemy assignments is thus an affine variety $P_2(\T)$, which is cut out by homogeneous quadratic polynomials. 

In this paper we define the Ptolemy field of a boundary-unipotent representation and show that it is isomorphic to the trace field. This gives rise to an efficient algorithm for \emph{exact} computation of the trace field of a hyperbolic manifold.

%Recall that the topological ideal triangulation $\T$ is a triangulation of the space $\widehat M$ obtained from $M$ by collapsing each boundary component to a point. It induces a triangulation of $M$ by truncated simplices. Let $P=\left\{\left(\begin{smallmatrix}1&x\\0&1\end{smallmatrix}\right)\right)\}$.

%The Ptolemy equations are homogeneous quadratic equations, and 
%there is an action of $x \in (\BC^*)^r$ on $c$ (where $r$ is the number of 
%boundary components of $M$) given by $(x \cdot c)_e=x' x'' c_e$ where 
%$x',x''$ are the coordinates of $x$ corresponding to the two ends of $e$.
 
%Let $V(\CT)$ denote the \emph{Ptolemy variety} and $V^\rd(\CT)$ denote the
%quotient of $V(\CT)$ by the action of $(\BC^*)^r$.  

\subsection{Decorated $\SL(2,\BC)$-representations}
\label{sub.decorated}
%A Ptolemy assignment explicitly determines (up to conjugation) a boundary-parabolic representation of $\pi_1(M)$ in $\SL(2,\BC)$. 
The precise relationship between Ptolemy assignments and representations is given by
\begin{equation}\label{eq:summaryof1to1}
\xymatrix{\left\{\txt{Points in\\$P_2(\T)$}\right\}\ar@{<->}[r]^-{1-1}&\left\{\txt{Natural $(\SL(2,\C),P)$-\\cocycles on $M$}\right\}\ar@{<->}[r]^-{1-1}&\left\{\txt{Generically decorated\\
$(\SL(2,\C),P)$-representations}\right\}.}
\end{equation}
The concepts are briefly described below, and the correspondences are illustrated in Figures~\ref{fig:DecToPtolemy} and~\ref{fig:PtolemyToCocycle}. We refer to Section~\ref{sec:Notation} for a summary of our notation.
The bijections of~\eqref{eq:summaryof1to1} first appeared in Zickert~\cite{ZickertDuke} (in a slightly different form), and were generalized to $\SL(n,\C)$-representations in Garoufalidis--Thurston--Zickert~\cite{GaroufalidisThurstonZickert}. 

\begin{itemize}[itemsep=1mm]
%\item \textit{Ptolemy assignment:} Assignment of non-zero complex numbers $c_{ij}$ to the oriented edges $\varepsilon_{ij}$ of each simplex of $\T$. Identified edges are assigned the same Ptolemy coordinate. For each simplex the Ptolemy relation~\eqref{eq:PtolemyRelation}.
%must be satisfied.
\item \textit{Natural cocycle:} Labeling of the edges of each truncated simplex by elements in $\SL(2,\C)$ satisfying the cocycle condition (the product around each face is $1$). The long edges are counter diagonal, i.e.~of the form $\left(\begin{smallmatrix}0&-x^{-1}\\x&0\end{smallmatrix}\right)$, and the short edges are \emph{non-trivial} elements in $P$. Identified edges are labeled by the same group element.
\item \textit{Decorated representation:} A \emph{decoration} of a boundary-parabolic representation $\rho$ is an assignment of a coset $gP$ to each vertex of $\Mtildehat$ which is equivariant with respect to $\rho$. A decoration is \emph{generic} if for each edge joining two vertices, the two $P$-cosets $gP$, $hP$ are distinct as $B$-cosets. This condition is equivalent to $\det(ge_1,he_1)\neq 0$. Two decorations are considered equal if they differ by left multiplication by a group element $g$.
\end{itemize}

%$\scalerel*{\widetilde{\mathcal{F}_K}}{j^2}$


%\begin{figure}[htb]
%\begin{center}
%\hspace{-10mm}
%\begin{minipage}[c]{0.3\textwidth}
%\scalebox{0.75}{\begingroup
 \setlength{\unitlength}{0.8pt}
 \begin{picture}(327.29865,196.42032)
 \put(0,0){\includegraphics{figures_gen/BasicPtolemy_included}}

\definecolor{inkcol1}{rgb}{0.0,0.0,0.0}
   \put(72.834381,14.946098){\rotatebox{360.0}{\makebox(0,0)[tl]{\strut{}{
    \begin{minipage}[h]{129.560208pt}
\textcolor{inkcol1}{\Large{$c_{02}$}}\\
\end{minipage}}}}}%


\definecolor{inkcol1}{rgb}{0.0,0.0,0.0}
   \put(200.606421,122.756458){\rotatebox{360.0}{\makebox(0,0)[tl]{\strut{}{
    \begin{minipage}[h]{129.560208pt}
\textcolor{inkcol1}{\Large{$c_{23}$}}\\
\end{minipage}}}}}%


\definecolor{inkcol1}{rgb}{0.0,0.0,0.0}
   \put(22.034991,114.899308){\rotatebox{360.0}{\makebox(0,0)[tl]{\strut{}{
    \begin{minipage}[h]{129.560208pt}
\textcolor{inkcol1}{\Large{$c_{01}$}}\\
\end{minipage}}}}}%


\definecolor{inkcol1}{rgb}{0.0,0.0,0.0}
   \put(126.320711,107.756458){\rotatebox{360.0}{\makebox(0,0)[tl]{\strut{}{
    \begin{minipage}[h]{129.560208pt}
\textcolor{inkcol1}{\Large{$c_{12}$}}\\
\end{minipage}}}}}%


\definecolor{inkcol1}{rgb}{0.0,0.0,0.0}
   \put(109.177841,133.470748){\rotatebox{360.0}{\makebox(0,0)[tl]{\strut{}{
    \begin{minipage}[h]{129.560208pt}
\textcolor{inkcol1}{\Large{$c_{03}$}}\\
\end{minipage}}}}}%


\definecolor{inkcol1}{rgb}{0.0,0.0,0.0}
   \put(133.463561,192.042178){\rotatebox{360.0}{\makebox(0,0)[tl]{\strut{}{
    \begin{minipage}[h]{129.560208pt}
\textcolor{inkcol1}{\Large{$c_{13}$}}\\
\end{minipage}}}}}%

 \end{picture}
\endgroup
}
%\end{minipage}
%\hspace{2mm}
%\begin{minipage}[c]{0.3\textwidth}
%\scalebox{0.75}{\begingroup
 \setlength{\unitlength}{0.8pt}
 \begin{picture}(234.20671,190.59085)
 \put(0,0){\includegraphics{figures_gen/TruncatedCocycle_included}}

\definecolor{inkcol1}{rgb}{0.0,0.0,0.0}
   \put(134.850526,94.132236){\rotatebox{360.0}{\makebox(0,0)[tl]{\strut{}{
    \begin{minipage}[h]{127.655896pt}
\textcolor{inkcol1}{\large{$\alpha$}}\\
\end{minipage}}}}}%


\definecolor{inkcol1}{rgb}{0.0,0.0,0.0}
   \put(80.534436,115.037126){\rotatebox{360.0}{\makebox(0,0)[tl]{\strut{}{
    \begin{minipage}[h]{127.655896pt}
\textcolor{inkcol1}{\large{$\beta$}}\\
\end{minipage}}}}}%

 \end{picture}
\endgroup
}
%\end{minipage}
%\hspace{2mm}
%\begin{minipage}[c]{0.3\textwidth}
%\scalebox{0.76}{\begingroup
 \setlength{\unitlength}{0.8pt}
 \begin{picture}(314.1539,200.47701)
 \put(0,0){\includegraphics{figures_gen/Decoration_included}}

\definecolor{inkcol1}{rgb}{0.0,0.0,0.0}
   \put(83.058597,177.913786){\rotatebox{360.0}{\makebox(0,0)[tl]{\strut{}{
    \begin{minipage}[h]{127.655896pt}
\textcolor{inkcol1}{\large{$g_1P$}}\\
\end{minipage}}}}}%


\definecolor{inkcol1}{rgb}{0.0,0.0,0.0}
   \put(-0.625263,14.222236){\rotatebox{360.0}{\makebox(0,0)[tl]{\strut{}{
    \begin{minipage}[h]{127.655896pt}
\textcolor{inkcol1}{\large{$g_0P$}}\\
\end{minipage}}}}}%


\definecolor{inkcol1}{rgb}{0.0,0.0,0.0}
   \put(184.302827,16.357046){\rotatebox{360.0}{\makebox(0,0)[tl]{\strut{}{
    \begin{minipage}[h]{127.655896pt}
\textcolor{inkcol1}{\large{$g_2P$}}\\
\end{minipage}}}}}%


\definecolor{inkcol1}{rgb}{0.0,0.0,0.0}
   \put(228.333317,203.419946){\rotatebox{360.0}{\makebox(0,0)[tl]{\strut{}{
    \begin{minipage}[h]{127.655896pt}
\textcolor{inkcol1}{\large{$g_3P$}}\\
\end{minipage}}}}}%

 \end{picture}
\endgroup
}
%\end{minipage}
%\hfill
%\\
%\begin{minipage}[t]{0.32\textwidth}
%\captionsetup{width=1\textwidth}
%\caption{Ptolemy assignment; the Ptolemy relation~\eqref{eq:PtolemyRelation} holds.}\label{fig:Ptolemy}
%\end{minipage}
%\hspace{-2mm}
%\begin{minipage}[t]{0.32\textwidth}
%\captionsetup{width=\textwidth}
%\caption[Bla]{Natural cocycle; $\alpha$ is counter-diagonal, $\beta\in P$}\label{fig:Cocycle}
%\end{minipage}
%\hspace{2mm}
%\begin{minipage}[t]{0.32\textwidth}
%\captionsetup{width=\textwidth}
%\caption[Bla]{Decoration; equivariant assignment of cosets.}\label{fig:Decoration}
%\end{minipage}
%\end{center}
%\end{figure}


%\begin{figure}[htb]
%\begin{center}
%\hspace{-10mm}
%\begin{minipage}[c]{0.48\textwidth}
%\scalebox{0.67}{\begingroup
 \setlength{\unitlength}{0.8pt}
 \begin{picture}(356.2648,197.84821)
 \put(0,0){\includegraphics{figures_gen/DecorationToPtolemy_included}}

\definecolor{inkcol1}{rgb}{0.0,0.0,0.0}
   \put(94.080375,203.534216){\rotatebox{360.0}{\makebox(0,0)[tl]{\strut{}{
    \begin{minipage}[h]{127.655896pt}
\textcolor{inkcol1}{\Large{$g_jP$}}\\
\end{minipage}}}}}%


\definecolor{inkcol1}{rgb}{0.0,0.0,0.0}
   \put(-0.205343,19.248506){\rotatebox{360.0}{\makebox(0,0)[tl]{\strut{}{
    \begin{minipage}[h]{127.655896pt}
\textcolor{inkcol1}{\Large{$g_iP$}}\\
\end{minipage}}}}}%


\definecolor{inkcol1}{rgb}{0.0,0.0,0.0}
   \put(158.653195,23.133176){\rotatebox{360.0}{\makebox(0,0)[tl]{\strut{}{
    \begin{minipage}[h]{127.655896pt}
\textcolor{inkcol1}{\Large{$i$}}\\
\end{minipage}}}}}%


\definecolor{inkcol1}{rgb}{0.0,0.0,0.0}
   \put(254.367485,202.418896){\rotatebox{360.0}{\makebox(0,0)[tl]{\strut{}{
    \begin{minipage}[h]{127.655896pt}
\textcolor{inkcol1}{\Large{$j$}}\\
\end{minipage}}}}}%


\definecolor{inkcol1}{rgb}{0.0,0.0,0.0}
   \put(216.429985,102.878696){\rotatebox{360.0}{\makebox(0,0)[tl]{\strut{}{
    \begin{minipage}[h]{127.655896pt}
\textcolor{inkcol1}{\Large{${c_{ij}=\det(g_ie_1,g_je_1)}$}}\\
\end{minipage}}}}}%

 \end{picture}
\endgroup
}
%\end{minipage}
%\hspace{2mm}
%\begin{minipage}[c]{0.48\textwidth}
%\scalebox{0.67}{\begingroup
 \setlength{\unitlength}{0.8pt}
 \begin{picture}(466.72302,188.14412)
 \put(0,0){\includegraphics{figures_gen/PtolemyToCocycleShort_included}}

\definecolor{inkcol1}{rgb}{0.0,0.0,0.0}
   \put(-4.751029,22.196066){\rotatebox{360.0}{\makebox(0,0)[tl]{\strut{}{
    \begin{minipage}[h]{127.655896pt}
\textcolor{inkcol1}{\Large{$i$}}\\
\end{minipage}}}}}%


\definecolor{inkcol1}{rgb}{0.0,0.0,0.0}
   \put(89.68401,195.207626){\rotatebox{360.0}{\makebox(0,0)[tl]{\strut{}{
    \begin{minipage}[h]{127.655896pt}
\textcolor{inkcol1}{\Large{$j$}}\\
\end{minipage}}}}}%


\definecolor{inkcol1}{rgb}{0.0,0.0,0.0}
   \put(180.09267,21.966456){\rotatebox{360.0}{\makebox(0,0)[tl]{\strut{}{
    \begin{minipage}[h]{127.655896pt}
\textcolor{inkcol1}{\Large{$k$}}\\
\end{minipage}}}}}%


\definecolor{inkcol1}{rgb}{0.0,0.0,0.0}
   \put(28.0647,120.024306){\rotatebox{360.0}{\makebox(0,0)[tl]{\strut{}{
    \begin{minipage}[h]{127.655896pt}
\textcolor{inkcol1}{\Large{$c_{ij}$}}\\
\end{minipage}}}}}%


\definecolor{inkcol1}{rgb}{0.0,0.0,0.0}
   \put(131.30331,119.758676){\rotatebox{360.0}{\makebox(0,0)[tl]{\strut{}{
    \begin{minipage}[h]{127.655896pt}
\textcolor{inkcol1}{\Large{$c_{jk}$}}\\
\end{minipage}}}}}%


\definecolor{inkcol1}{rgb}{0.0,0.0,0.0}
   \put(74.22969,14.702816){\rotatebox{360.0}{\makebox(0,0)[tl]{\strut{}{
    \begin{minipage}[h]{127.655896pt}
\textcolor{inkcol1}{\Large{$c_{ki}$}}\\
\end{minipage}}}}}%


\definecolor{inkcol1}{rgb}{0.0,0.0,0.0}
   \put(206.12483,21.323636){\rotatebox{360.0}{\makebox(0,0)[tl]{\strut{}{
    \begin{minipage}[h]{127.655896pt}
\textcolor{inkcol1}{\Large{$i$}}\\
\end{minipage}}}}}%


\definecolor{inkcol1}{rgb}{0.0,0.0,0.0}
   \put(304.60049,193.325046){\rotatebox{360.0}{\makebox(0,0)[tl]{\strut{}{
    \begin{minipage}[h]{127.655896pt}
\textcolor{inkcol1}{\Large{$j$}}\\
\end{minipage}}}}}%


\definecolor{inkcol1}{rgb}{0.0,0.0,0.0}
   \put(396.77039,22.471536){\rotatebox{360.0}{\makebox(0,0)[tl]{\strut{}{
    \begin{minipage}[h]{127.655896pt}
\textcolor{inkcol1}{\Large{$k$}}\\
\end{minipage}}}}}%


\definecolor{inkcol1}{rgb}{0.0,0.0,0.0}
   \put(305.04954,117.233296){\rotatebox{360.0}{\makebox(0,0)[tl]{\strut{}{
    \begin{minipage}[h]{127.655896pt}
\textcolor{inkcol1}{\LARGE{${\beta\!=\!\left(\begin{smallmatrix}1&-\frac{c_{ki}}{c_{ij}c_{jk}}\\0&1\end{smallmatrix}\right)}$}}\\
\end{minipage}}}}}%


\definecolor{inkcol1}{rgb}{0.0,0.0,0.0}
   \put(135.5529,177.711596){\rotatebox{360.0}{\makebox(0,0)[tl]{\strut{}{
    \begin{minipage}[h]{127.655896pt}
\textcolor{inkcol1}{\LARGE{${\alpha\!=\!\left(\begin{smallmatrix}0&-c_{ij}^{-1}\\c_{ij}&0\end{smallmatrix}\right)}$}}\\
\end{minipage}}}}}%


\definecolor{inkcol1}{rgb}{0.0,0.0,0.0}
   \put(262.66347,88.343146){\rotatebox{360.0}{\makebox(0,0)[tl]{\strut{}{
    \begin{minipage}[h]{127.655896pt}
\textcolor{inkcol1}{\Large{$\alpha$}}\\
\end{minipage}}}}}%


\definecolor{inkcol1}{rgb}{0.0,0.0,0.0}
   \put(300.24217,130.696496){\rotatebox{360.0}{\makebox(0,0)[tl]{\strut{}{
    \begin{minipage}[h]{127.655896pt}
\textcolor{inkcol1}{\Large{$\beta$}}\\
\end{minipage}}}}}%

 \end{picture}
\endgroup
}
%\end{minipage}
%\hfill
%\\
%\begin{minipage}[t]{0.48\textwidth}
%\captionsetup{width=1\textwidth}
%\caption{From decorations to Ptolemy assignments.}\label{fig:DecToPtolemy}
%\end{minipage}
%\hspace{-2mm}
%\begin{minipage}[t]{0.48\textwidth}
%\captionsetup{width=\textwidth}
%\caption[Bla]{From Ptolemy assignments to natural cocycles.}\label{fig:PtolemyToCocycle}
%\end{minipage}
%\hspace{2mm}
%\end{center}
%\end{figure}



\begin{figure}[htb]
\begin{center}
\hspace{-30mm}
\begin{minipage}[c]{0.3\textwidth}
\scalebox{0.65}{\begingroup
 \setlength{\unitlength}{0.8pt}
 \begin{picture}(327.29865,196.42032)
 \put(0,0){\includegraphics{figures_gen/BasicPtolemy_included}}

\definecolor{inkcol1}{rgb}{0.0,0.0,0.0}
   \put(72.834381,14.946098){\rotatebox{360.0}{\makebox(0,0)[tl]{\strut{}{
    \begin{minipage}[h]{129.560208pt}
\textcolor{inkcol1}{\Large{$c_{02}$}}\\
\end{minipage}}}}}%


\definecolor{inkcol1}{rgb}{0.0,0.0,0.0}
   \put(200.606421,122.756458){\rotatebox{360.0}{\makebox(0,0)[tl]{\strut{}{
    \begin{minipage}[h]{129.560208pt}
\textcolor{inkcol1}{\Large{$c_{23}$}}\\
\end{minipage}}}}}%


\definecolor{inkcol1}{rgb}{0.0,0.0,0.0}
   \put(22.034991,114.899308){\rotatebox{360.0}{\makebox(0,0)[tl]{\strut{}{
    \begin{minipage}[h]{129.560208pt}
\textcolor{inkcol1}{\Large{$c_{01}$}}\\
\end{minipage}}}}}%


\definecolor{inkcol1}{rgb}{0.0,0.0,0.0}
   \put(126.320711,107.756458){\rotatebox{360.0}{\makebox(0,0)[tl]{\strut{}{
    \begin{minipage}[h]{129.560208pt}
\textcolor{inkcol1}{\Large{$c_{12}$}}\\
\end{minipage}}}}}%


\definecolor{inkcol1}{rgb}{0.0,0.0,0.0}
   \put(109.177841,133.470748){\rotatebox{360.0}{\makebox(0,0)[tl]{\strut{}{
    \begin{minipage}[h]{129.560208pt}
\textcolor{inkcol1}{\Large{$c_{03}$}}\\
\end{minipage}}}}}%


\definecolor{inkcol1}{rgb}{0.0,0.0,0.0}
   \put(133.463561,192.042178){\rotatebox{360.0}{\makebox(0,0)[tl]{\strut{}{
    \begin{minipage}[h]{129.560208pt}
\textcolor{inkcol1}{\Large{$c_{13}$}}\\
\end{minipage}}}}}%

 \end{picture}
\endgroup
}
\end{minipage}
\hspace{-10mm}
\begin{minipage}[c]{0.3\textwidth}
\scalebox{0.65}{\begingroup
 \setlength{\unitlength}{0.8pt}
 \begin{picture}(234.20671,190.59085)
 \put(0,0){\includegraphics{figures_gen/TruncatedCocycle_included}}

\definecolor{inkcol1}{rgb}{0.0,0.0,0.0}
   \put(134.850526,94.132236){\rotatebox{360.0}{\makebox(0,0)[tl]{\strut{}{
    \begin{minipage}[h]{127.655896pt}
\textcolor{inkcol1}{\large{$\alpha$}}\\
\end{minipage}}}}}%


\definecolor{inkcol1}{rgb}{0.0,0.0,0.0}
   \put(80.534436,115.037126){\rotatebox{360.0}{\makebox(0,0)[tl]{\strut{}{
    \begin{minipage}[h]{127.655896pt}
\textcolor{inkcol1}{\large{$\beta$}}\\
\end{minipage}}}}}%

 \end{picture}
\endgroup
}
\end{minipage}
%\hspace{2mm}
\begin{minipage}[c]{0.3\textwidth}
\scalebox{0.62}{\begingroup
 \setlength{\unitlength}{0.8pt}
 \begin{picture}(466.72302,188.14412)
 \put(0,0){\includegraphics{figures_gen/PtolemyToCocycleShort_included}}

\definecolor{inkcol1}{rgb}{0.0,0.0,0.0}
   \put(-4.751029,22.196066){\rotatebox{360.0}{\makebox(0,0)[tl]{\strut{}{
    \begin{minipage}[h]{127.655896pt}
\textcolor{inkcol1}{\Large{$i$}}\\
\end{minipage}}}}}%


\definecolor{inkcol1}{rgb}{0.0,0.0,0.0}
   \put(89.68401,195.207626){\rotatebox{360.0}{\makebox(0,0)[tl]{\strut{}{
    \begin{minipage}[h]{127.655896pt}
\textcolor{inkcol1}{\Large{$j$}}\\
\end{minipage}}}}}%


\definecolor{inkcol1}{rgb}{0.0,0.0,0.0}
   \put(180.09267,21.966456){\rotatebox{360.0}{\makebox(0,0)[tl]{\strut{}{
    \begin{minipage}[h]{127.655896pt}
\textcolor{inkcol1}{\Large{$k$}}\\
\end{minipage}}}}}%


\definecolor{inkcol1}{rgb}{0.0,0.0,0.0}
   \put(28.0647,120.024306){\rotatebox{360.0}{\makebox(0,0)[tl]{\strut{}{
    \begin{minipage}[h]{127.655896pt}
\textcolor{inkcol1}{\Large{$c_{ij}$}}\\
\end{minipage}}}}}%


\definecolor{inkcol1}{rgb}{0.0,0.0,0.0}
   \put(131.30331,119.758676){\rotatebox{360.0}{\makebox(0,0)[tl]{\strut{}{
    \begin{minipage}[h]{127.655896pt}
\textcolor{inkcol1}{\Large{$c_{jk}$}}\\
\end{minipage}}}}}%


\definecolor{inkcol1}{rgb}{0.0,0.0,0.0}
   \put(74.22969,14.702816){\rotatebox{360.0}{\makebox(0,0)[tl]{\strut{}{
    \begin{minipage}[h]{127.655896pt}
\textcolor{inkcol1}{\Large{$c_{ki}$}}\\
\end{minipage}}}}}%


\definecolor{inkcol1}{rgb}{0.0,0.0,0.0}
   \put(206.12483,21.323636){\rotatebox{360.0}{\makebox(0,0)[tl]{\strut{}{
    \begin{minipage}[h]{127.655896pt}
\textcolor{inkcol1}{\Large{$i$}}\\
\end{minipage}}}}}%


\definecolor{inkcol1}{rgb}{0.0,0.0,0.0}
   \put(304.60049,193.325046){\rotatebox{360.0}{\makebox(0,0)[tl]{\strut{}{
    \begin{minipage}[h]{127.655896pt}
\textcolor{inkcol1}{\Large{$j$}}\\
\end{minipage}}}}}%


\definecolor{inkcol1}{rgb}{0.0,0.0,0.0}
   \put(396.77039,22.471536){\rotatebox{360.0}{\makebox(0,0)[tl]{\strut{}{
    \begin{minipage}[h]{127.655896pt}
\textcolor{inkcol1}{\Large{$k$}}\\
\end{minipage}}}}}%


\definecolor{inkcol1}{rgb}{0.0,0.0,0.0}
   \put(305.04954,117.233296){\rotatebox{360.0}{\makebox(0,0)[tl]{\strut{}{
    \begin{minipage}[h]{127.655896pt}
\textcolor{inkcol1}{\LARGE{${\beta\!=\!\left(\begin{smallmatrix}1&-\frac{c_{ki}}{c_{ij}c_{jk}}\\0&1\end{smallmatrix}\right)}$}}\\
\end{minipage}}}}}%


\definecolor{inkcol1}{rgb}{0.0,0.0,0.0}
   \put(135.5529,177.711596){\rotatebox{360.0}{\makebox(0,0)[tl]{\strut{}{
    \begin{minipage}[h]{127.655896pt}
\textcolor{inkcol1}{\LARGE{${\alpha\!=\!\left(\begin{smallmatrix}0&-c_{ij}^{-1}\\c_{ij}&0\end{smallmatrix}\right)}$}}\\
\end{minipage}}}}}%


\definecolor{inkcol1}{rgb}{0.0,0.0,0.0}
   \put(262.66347,88.343146){\rotatebox{360.0}{\makebox(0,0)[tl]{\strut{}{
    \begin{minipage}[h]{127.655896pt}
\textcolor{inkcol1}{\Large{$\alpha$}}\\
\end{minipage}}}}}%


\definecolor{inkcol1}{rgb}{0.0,0.0,0.0}
   \put(300.24217,130.696496){\rotatebox{360.0}{\makebox(0,0)[tl]{\strut{}{
    \begin{minipage}[h]{127.655896pt}
\textcolor{inkcol1}{\Large{$\beta$}}\\
\end{minipage}}}}}%

 \end{picture}
\endgroup
}
\end{minipage}
\\
\hspace{-20mm}
\begin{minipage}[t]{0.32\textwidth}
\captionsetup{width=1\textwidth}
\caption{Ptolemy assignment; the Ptolemy relation~\eqref{eq:PtolemyRelation} holds.}\label{fig:Ptolemy}
\end{minipage}
\hspace{-10mm}
\begin{minipage}[t]{0.32\textwidth}
\captionsetup{width=\textwidth}
\caption[Bla]{Natural cocycle; $\alpha$ is counter-diagonal, $\beta\in P$.}\label{fig:Cocycle}
\end{minipage}
\hspace{5mm}
\begin{minipage}[t]{0.32\textwidth}
\captionsetup{width=\textwidth}
\caption[Bla]{From Ptolemy assignments to natural cocycles.}\label{fig:PtolemyToCocycle}
\end{minipage}
\end{center}
\end{figure}


\begin{figure}[htb]
\begin{center}
\hspace{4mm}
\begin{minipage}[c]{0.3\textwidth}
\scalebox{0.7}{\begingroup
 \setlength{\unitlength}{0.8pt}
 \begin{picture}(314.1539,200.47701)
 \put(0,0){\includegraphics{figures_gen/Decoration_included}}

\definecolor{inkcol1}{rgb}{0.0,0.0,0.0}
   \put(83.058597,177.913786){\rotatebox{360.0}{\makebox(0,0)[tl]{\strut{}{
    \begin{minipage}[h]{127.655896pt}
\textcolor{inkcol1}{\large{$g_1P$}}\\
\end{minipage}}}}}%


\definecolor{inkcol1}{rgb}{0.0,0.0,0.0}
   \put(-0.625263,14.222236){\rotatebox{360.0}{\makebox(0,0)[tl]{\strut{}{
    \begin{minipage}[h]{127.655896pt}
\textcolor{inkcol1}{\large{$g_0P$}}\\
\end{minipage}}}}}%


\definecolor{inkcol1}{rgb}{0.0,0.0,0.0}
   \put(184.302827,16.357046){\rotatebox{360.0}{\makebox(0,0)[tl]{\strut{}{
    \begin{minipage}[h]{127.655896pt}
\textcolor{inkcol1}{\large{$g_2P$}}\\
\end{minipage}}}}}%


\definecolor{inkcol1}{rgb}{0.0,0.0,0.0}
   \put(228.333317,203.419946){\rotatebox{360.0}{\makebox(0,0)[tl]{\strut{}{
    \begin{minipage}[h]{127.655896pt}
\textcolor{inkcol1}{\large{$g_3P$}}\\
\end{minipage}}}}}%

 \end{picture}
\endgroup
}
\end{minipage}
\hfill
\begin{minipage}[c]{0.48\textwidth}
\scalebox{0.63}{\begingroup
 \setlength{\unitlength}{0.8pt}
 \begin{picture}(356.2648,197.84821)
 \put(0,0){\includegraphics{figures_gen/DecorationToPtolemy_included}}

\definecolor{inkcol1}{rgb}{0.0,0.0,0.0}
   \put(94.080375,203.534216){\rotatebox{360.0}{\makebox(0,0)[tl]{\strut{}{
    \begin{minipage}[h]{127.655896pt}
\textcolor{inkcol1}{\Large{$g_jP$}}\\
\end{minipage}}}}}%


\definecolor{inkcol1}{rgb}{0.0,0.0,0.0}
   \put(-0.205343,19.248506){\rotatebox{360.0}{\makebox(0,0)[tl]{\strut{}{
    \begin{minipage}[h]{127.655896pt}
\textcolor{inkcol1}{\Large{$g_iP$}}\\
\end{minipage}}}}}%


\definecolor{inkcol1}{rgb}{0.0,0.0,0.0}
   \put(158.653195,23.133176){\rotatebox{360.0}{\makebox(0,0)[tl]{\strut{}{
    \begin{minipage}[h]{127.655896pt}
\textcolor{inkcol1}{\Large{$i$}}\\
\end{minipage}}}}}%


\definecolor{inkcol1}{rgb}{0.0,0.0,0.0}
   \put(254.367485,202.418896){\rotatebox{360.0}{\makebox(0,0)[tl]{\strut{}{
    \begin{minipage}[h]{127.655896pt}
\textcolor{inkcol1}{\Large{$j$}}\\
\end{minipage}}}}}%


\definecolor{inkcol1}{rgb}{0.0,0.0,0.0}
   \put(216.429985,102.878696){\rotatebox{360.0}{\makebox(0,0)[tl]{\strut{}{
    \begin{minipage}[h]{127.655896pt}
\textcolor{inkcol1}{\Large{${c_{ij}=\det(g_ie_1,g_je_1)}$}}\\
\end{minipage}}}}}%

 \end{picture}
\endgroup
}
\end{minipage}
\hfill
\\
\begin{minipage}[t]{0.48\textwidth}
\captionsetup{width=1\textwidth}
\caption{Decoration; equivariant assignment of cosets.}\label{fig:Decoration}
\end{minipage}
\hspace{-2mm}
\begin{minipage}[t]{0.48\textwidth}
\captionsetup{width=\textwidth}
\caption[Bla]{From decorations to Ptolemy assignments.}\label{fig:DecToPtolemy}
\end{minipage}
%\hspace{2mm}
\end{center}
\end{figure}



By ignoring the decoration, \eqref{eq:summaryof1to1} yields a map
\begin{equation}\label{eq:PtolemyToRep}
\mathcal R\colon P_2(\T)\to \big\{(\SL(2,\C),P)\text{-representations}\big\}
\big/\Conj.
\end{equation}
The representation corresponding to a Ptolemy assignment is given explicitly in terms of the natural cocycle.

\begin{remark}
Note that a natural cocycle canonically determines a representation of the edge path groupoid of the triangulation of $M$ by truncated simplices.
\end{remark}

\begin{remark} 
Every boundary-parabolic representation has a decoration, but a representation may have only non-generic decorations. %This happens, e.g., when the triangulation has a non-essential edge (see \cite{GaroufalidisThurstonZickert}). 
The map $\mathcal R$ is thus not surjective in general, and the image depends on the triangulation. However, if the triangulation is sufficiently fine, 
$\mathcal R$ is surjective (see~\cite{GaroufalidisThurstonZickert}). The preimage of a representation depends on the image of the peripheral subgroups (see Proposition~\ref{prop:GenericBoundaryNonTrivial}).\end{remark}



\subsection{Obstruction classes and $\PSL(2,\C)$-representations}
There is a subtle distinction between representations in $\SL(2,\C)$ versus $\PSL(2,\C)$. The geometric representation of a hyperbolic manifold always lifts to an $\SL(2,\C)$-representation, but for a one-cusped manifold, no lift is boundary-parabolic (any lift will take a longitude to an element of trace $-2$~\cite{Calegari}). %In fact, some manifolds have no interesting boundary-parabolic $\SL(2,\C)$-representations. The figure-$8$ knot, for example, has only abelian such.

The obstruction to lifting a boundary-parabolic $\PSL(2,\C)$-representation to a boundary-parabolic $\SL(2,\C)$-representation is a class in $H^2(\widehat M;\Z/2\Z)$. For each such class there is a Ptolemy variety $P^\sigma_2(\T)$, which maps to the set of $\PSL(2,\C)$-representations with obstruction class $\sigma$. More precisely, $P^\sigma_2(\T)$ is defined for each $2$-cocycle $\sigma\in Z^2(\widehat M;\Z/2\Z)$, and up to canonical isomorphism only depends on the cohomology class of $\sigma$. The Ptolemy variety for the trivial cocycle equals $P_2(\T)$. The analogue of~\eqref{eq:summaryof1to1} is
\begin{equation}\label{eq:1to1Obstruction}
\xymatrix{\left\{\txt{Points in\\$P^{\sigma}_2(\T)$}\right\}\ar@{<->}[r]^-{1-1}&\left\{\txt{Lifted natural\\$(\PSL(2,\C),P)$-cocycles\\with obstruction cocycle $\sigma$}\right\}\ar@{->>}[r]&\left\{\txt{Generically decorated\\
$(\PSL(2,\C),P)$-representations\\with obstruction class $\sigma$}\right\}.}
\end{equation}
A lifted natural cocycle is defined as above, except that the product along a face is now $\pm I$, where the sign is determined by $\sigma$. The right map is no longer a $1$-$1$ correspondence; the preimage of each decorated representation is the choice of lifts, i.e.~parametrized by a cocycle in $Z^1(\widehat M;\Z/2\Z)$. We refer to \cite{GaroufalidisThurstonZickert} for details. As in~\eqref{eq:PtolemyToRep}, ignoring the decoration yields a map
\begin{equation}\label{eq:PtolemyToRepObstruction}
\mathcal R\colon P_2^{\sigma}(\T)\to \left\{\txt{$(\PSL(2,\C),P)$-representations\\with obstruction class $\sigma$}\right\}
\big/\Conj.
\end{equation}
which is explictly given in terms of the natural cocycle.
\begin{theorem}[Garoufalidis--Thurston--Zickert~\cite{GaroufalidisThurstonZickert}]\label{thm:EssentialEdges}
If $M$ is hyperbolic, and all edges of $\T$ are essential, the geometric representation is in the image of $\mathcal R$.
\end{theorem}
\begin{remark}
If $\T$ has a non-essential edge, all Ptolemy varieties will be empty. Hence, if $P^\sigma_2(\T)$ is non-empty for some $\sigma$, and if $M$ is hyperbolic, the geometric representation is detected by the Ptolemy variety of the geometric obstruction class.
\end{remark}


\subsection{Our results}
We view the Ptolemy varieties $P^\sigma_2(\T)$ as subsets of an ambient space $\C^e$, with coordinates indexed by the $1$-cells of $\T$. Let $T=(\C^*)^v$, with the coordinates indexed by the boundary components of $M$. 

\begin{definition} The \emph{diagonal action} is the action of $T$ on $P^\sigma_2(\T)$ where $(x_1,\dots,x_v)\in T$ acts on a Ptolemy assignment by replacing the Ptolemy coordinate $c$ of an edge $e$ with $x_ix_jc$, where $x_i$ and $x_j$ are the coordinates corresponding to the ends of $e$. Let 
\begin{equation}
P^\sigma_2(\T)_\rd=P^\sigma_2(\T)\big/T.
\end{equation}
\end{definition}


\begin{definition}
A boundary-parabolic $\PSL(2,\C)$-representation is \emph{generic} if it has a generic decoration. It is \emph{boundary-non-trivial} if each peripheral subgroup has non-trivial image.
\end{definition}
\begin{remark}
Note that the notion of genericity is with respect to the triangulation. By Theorem~\ref{thm:EssentialEdges}, if all edges of $\T$ are essential (and $\T$ has no interior vertices), the geometric representation of a cusped hyperbolic manifold is always generic and boundary-non-trivial.
\end{remark}
\begin{remark}
Note that if $M$ has spherical boundary components (e.g.~if $\T$ is a triangulation of a closed manifold), no representation is boundary-non-trivial.
\end{remark}

\begin{proposition}\label{prop:GenericBoundaryNonTrivial}
The map $\mathcal R$ in~\eqref{eq:PtolemyToRepObstruction} factors through $P_2^{\sigma}(\T)_{\rd}$, i.e.,~we have
\begin{equation}\label{eq:GenericBoundaryNonTrivial}
\cxymatrix{{\mathcal R\colon P_2^{\sigma}(\T)_{\rd}\ar@{->}[r]& \left\{\txt{$(\PSL(2,\C),P)$-representations\\with obstruction class $\sigma$}\right\}\big/\Conj.}}
\end{equation}
The image is the set of generic representations, and the preimage of a generic, boundary-non-trivial representation is finite and parametrized by $H^1(\widehat M;\Z/2\Z)$. %Moreover, if $c\in P_2^{\sigma}(\T)_{\rd}$ is in a $0$-dimensional component, the image is boundary-non-trivial.
\end{proposition}
\begin{remark}
For the corresponding map from $P_2(\T)_\rd$ to $(\SL(2,\C),P)$-representations, the preimage of a generic boundary-non-trivial representation is a single point.
\end{remark}

\begin{remark}\label{rm:HigherComponents}
The preimage of a representation which is not boundary-non-trivial is never finite. In fact its dimension is the number of boundary-components that are collapsed. In particular, it follows that if $c\in P_2^{\sigma}(\T)_{\rd}$ is in a $0$-dimensional component (which is not contained in a higher dimensional component), the image is boundary-non-trivial.
\end{remark}

By geometric invariant theory, $P^\sigma_2(\T)_\rd$ is a variety whose coordinate ring is the ring of invariants $\mathcal O^T$ of the coordinate ring $\mathcal O$ of $P^\sigma_2(T)$.

\begin{definition}\label{def:PtolemyField}
Let $c\in P^\sigma_2(\T)$. The Ptolemy field of $c$ is the field
\begin{equation}
k_c=\Q\left(\big\{p(c_1,\dots,c_e)\mid p\in\mathcal O^T\big\}\right).
\end{equation}
The Ptolemy field of a generic boundary-non-trivial representation is the Ptolemy field of any preimage under~\eqref{eq:GenericBoundaryNonTrivial}.
\end{definition}
Clearly, the Ptolemy field only depends on the image in $P^\sigma_2(\T)_\rd$. Our main result is the following.
\begin{theorem}\label{thm:PtolemyFieldEqualsTraceField}
The Ptolemy field of a boundary-non-trivial, generic, boundary-parabolic representation $\rho$ in $\PSL(2,\C)$ or $\SL(2,\C)$ is equal to its trace field.\qed
\end{theorem}
\begin{remark}
For a cusped hyperbolic $3$-manifold the \emph{shape field} is in general smaller than the trace field. The shape field equals the \emph{invariant trace field} (see, e.g.,~\cite{MaclachlanReid}).
\end{remark}
For computations of the Ptolemy field, we need an explicit description of the ring of invariants $\mathcal O^T$, or equivalently, the reduced Ptolemy variety $P^\sigma_2(\T)_\rd$.
%The result below gives an embedding of $P_2^{\sigma}(\T)_{\rd}$ in $\C^e$, which is used in the definition of the Ptolemy field below.
\begin{proposition}\label{prop:FixDecoration}
There exist $1$-cells $\varepsilon_1,\dots,\varepsilon_v$ of $\T$ such that the reduced Ptolemy variety $P_2^{\sigma}(\T)_{\rd}$ is naturally isomorphic to the subvariety of $P^\sigma_2(\T)$ obtained by intersecting with the affine hyperplane $c_{\varepsilon_1}=\dots=c_{\varepsilon_v}=1$. % \qed
\end{proposition}
%Such $1$-cells are called~\emph{basic} (see Definition~\ref{}), and the corresponding embedding in $\C^e$ is called a \emph{basic embedding}.
\begin{corollary}
Let $c\in P^\sigma_2(\T)_\rd$. Under an isomorphism as in Proposition~\ref{prop:FixDecoration}, the Ptolemy field of $c$ is the field generated by the Ptolemy coordinates.\qed %In particular, if $c$ is in a zero-dimensional component, the trace field of the corresponding representation (which is boundary-non-trivial by Remark~\ref{rm:HigherComponents}) is the field generated by the Ptolemy coordinates under the embedding . \qed
\end{corollary}
\begin{remark} 
A concrete method for selecting $1$-cells as in Proposition~\ref{prop:FixDecoration} is described in Section~\ref{sec:AlphaStarBasic}.
\end{remark}
Analogues of our results for higher rank Ptolemy varieties are discussed in Section~\ref{sec:HigherPtolemy}. The analogue of Proposition~\ref{prop:GenericBoundaryNonTrivial} holds for representations that are \emph{boundary-non-degenerate} (see Definition~\ref{def:BoundaryNonDeg}), and the analogue of Proposition~\ref{prop:FixDecoration} leads to a simple computation of the Ptolemy field. 
\begin{conjecture}
The Ptolemy field of a boundary-non-degenerate, generic, boundary-unipotent representation $\rho$ in $\SL(n,\C)$ or $\PSL(n,\C)$ is equal to its trace field.\qed
\end{conjecture}
%For the higher rank Ptolemy varieties, we prove analogues of  and \ref{prop:FixDecoration} in  A higher rank representation is \emph{boundary-non-degenerate} if the normalizer of the image of each peripheral subgroup is isomorphic to the Borel subgroup $B$ of upper diagonal matrices.
\begin{remark}
The computation of reduced Ptolemy varieties is remarkably efficient using Magma~\cite{Magma}.
For all but a few census manifolds, primary decompositions of the (reduced) Ptolemy varieties $P^\sigma_2(\T)$ can be computed in a fraction of a second on a standard laptop. A database of representations can be found at \url{http://ptolemy.unhyperbolic.org/}~\cite{Unhyperbolic}.
All our tools have been incorporated into SnapPy~\cite{SnapPy} by the second author and the Ptolemy fields obtained through:

\begin{verbatim}
>>> from snappy import Manifold
>>> p=Manifold("m019").ptolemy_variety(2,'all')
>>> p.retrieve_solutions().number_field()
... [[x^4 - 2*x^2 - 3*x - 1], [x^4 + x - 1]]
 \end{verbatim}
 
The number fields are grouped by obstruction class. In this example, we see that the Ptolemy variety for the non-trivial obstruction class has a component with number field $x^4+x-1$, which is the trace field of m019. The above code retrieves a precomputed decomposition of the Ptolemy variety from~\cite{Unhyperbolic}.
In sage or SnapPy with Magma installed, you can use \texttt{p.compute\_solutions().number\_field()} to compute the decomposition.
\end{remark}


\section{Notation}\label{sec:Notation}
\subsection{Triangulations}
Let $M$ be a compact oriented $3$-manifold with (possibly empty) boundary. We refer to the boundary components as \emph{cusps} (although they may not be tori). Let $\widehat M$ be the space obtained from $M$ by collapsing each boundary component to a point.
\begin{definition}
A (concrete) \emph{triangulation} of $M$ is an identification of $\widehat M$ with a space obtained from a collection of simplices by gluing together pairs of faces by affine homeomorphisms. For each simplex $\Delta$ of $\T$ we fix an identification of $\Delta$ with a standard simplex.  
\end{definition}
\begin{remark}
By drilling out disjoint balls if necessary (this does not change the fundamental group), we may assume that the triangulation of $M$ is \emph{ideal}, i.e.~that each $0$-cell corresponds to a boundary-component of $M$. For example, we regard a triangulation of a closed manifold as an ideal triangulation of a manifold with boundary a union of spheres.
\end{remark}
%The identification with a standard simplex gives a convenient labeling scheme for the edges. Namely, we label the (oriented) edge from vertex $i$ to $j$ by $\varepsilon_{ij}$. 
\begin{definition}
%A triangulation is \emph{ordered} if $\varepsilon_{ij,k}\sim \varepsilon_{i^\prime j^\prime,k^\prime}$ implies $i<j\iff i^\prime<j^\prime$.
A triangulation is \emph{oriented} if the identifications with standard simplices are orientation preserving.
\end{definition}
\begin{remark}
All of the triangulations in the SnapPy censuses \texttt{OrientableCuspedCensus}, \texttt{LinkExteriors} and \texttt{HTLinkExteriors}~\cite{SnapPy} are oriented. Unless otherwise specified we shall assume that our triangulations are oriented.
\end{remark}


A triangulation gives rise to a triangulation of $M$ by truncated simplices, and to a triangulation of $\Mtildehat$.

\subsection{Miscellaneous}
\begin{itemize}
\item The number of vertices, edges, faces, and simplices, of a triangulation $\T$ are denoted by $v$, $e$, $f$, and $s$, respectively.
\item The standard basic vectors in $\Z^k$ are denoted by $e_1,\dots,e_{k}$. 
\item The (oriented) edge of simplex $k$ from vertex $i$ to $j$ is denoted by $\varepsilon_{ij,k}$. 
\item The matrix groups $\left\{\left(\begin{smallmatrix}1&x\\0&1\end{smallmatrix}\right)\right\}$ and $\left\{\left(\begin{smallmatrix}a&x\\0&a^{-1}\end{smallmatrix}\right)\right\}$ are denoted by $P$ and $B$, respectively. The higher rank analogue of $P$ is denoted by $N$.
\item A representation is \emph{boundary-parabolic} if it takes each peripheral subroup to a conjugate of $P$. Such is also called a \emph{$(G,P)$-representation} ($G=\SL(2,\C)$ or $\PSL(2,\C)$). In the higher rank case, such a representation is called boundary-unipotent.
\item A triangulation is \emph{ordered} if $\varepsilon_{ij,k}\sim \varepsilon_{i^\prime j^\prime,k^\prime}$ implies that $i<j\iff i^\prime<j^\prime$.
\end{itemize}


\section{The Ptolemy varieties}\label{sec:PtolemyVariety}
We define the Ptolemy variety for $n=2$ following Garoufalidis--Thurston--Zickert~\cite{GaroufalidisThurstonZickert} (see also Garoufalidis--Goerner--Zickert~\cite{GaroufalidisGoernerZickert}).
%Say something brief about obstruction classes
\subsection{The $\SL(2,\C)$-Ptolemy variety}
Assign to each oriented edge $\varepsilon_{ij,k}$ of $\Delta_k\in \T$ a \emph{Ptolemy coordinate $c_{ij,k}$}. Consider the affine algebraic set $A$ defined by the \emph{Ptolemy relations}
\begin{equation}\label{eq:PtolemyRelationijk}
c_{03,k}c_{12,k}+c_{01,k}c_{23,k}=c_{02,k}c_{13,k}, \qquad k=1,2,\dots,t,
\end{equation}
the \emph{identification relations}
\begin{equation}\label{eq:IdentificationRelations}
c_{ij,k}=c_{i^\prime j^\prime,k^\prime}\qquad \text{when}\qquad \varepsilon_{ij,k}\sim\varepsilon_{i^\prime j^\prime,k^\prime},
\end{equation}
and the \emph{edge orientation relations} $c_{ij,k}=-c_{ji,k}$. By only considering $i<j$, we shall always eliminate the edge orientation relations.
\begin{definition}\label{def:PtolemyVariety} The \emph{Ptolemy variety} $P_2(\T)$ is the Zariski open subset of $A$ consisting of points with non-zero Ptolemy coordinates.
\end{definition}
\begin{remark} One can concretely obtain $P_2(\T)$ from $A$ by adding a dummy variable $x$ and a dummy relation $x\prod c_{ij,k}=1$.
\end{remark}

\begin{remark}\label{rm:PickRepresentative} We can eliminate the identification relations~\eqref{eq:IdentificationRelations} by selecting a representative for each edge cycle. This gives an embedding of the Ptolemy variety in an ambient space $\C^e$, where it is cut out by $s$ Ptolemy relations, one for each simplex. Note that when all boundary components are tori, $s=e$.
\end{remark}

\subsubsection{The figure-$8$ knot}
Consider the ideal triangulation of the figure-$8$ knot complement shown in Figure~\ref{fig:Fig8Ordered}. The Ptolemy variety $P_2(\T)$ is given by
\begin{equation}\label{eq:Fig8PtolemyUnsimplified}
\begin{aligned}[c]
c_{03,0}c_{12,0}+c_{01,0}c_{23,0}&=c_{02,0}c_{13,0},\\
c_{03,1}c_{12,1}+c_{01,1}c_{23,1}&=c_{02,1}c_{13,1},
\end{aligned}
\qquad 
\begin{aligned}[c]
c_{02,0}&=c_{12,0}=c_{13,0}=c_{01,1}=c_{03,1}=c_{23,1}\\
c_{01,0}&=c_{03,0}=c_{23,0}=c_{02,1}=c_{12,1}=c_{13,1}.
\end{aligned}
\end{equation}
By selecting representatives $\varepsilon_{23,0}$ and $\varepsilon_{13,0}$ for the $2$ edge cycles, $P_2(\T)$ embeds in $\C^2$ where it is given by
\begin{equation}\label{eq:Fig8PtolemyTrivial}
c_{23,0}c_{13,0}+c_{23,0}^2=c_{13,0}^2,\qquad c_{13,0}c_{23,0}+c_{13,0}^2=c_{23,0}^2.
\end{equation}
It follows that $P_2(\T)$ is empty, which is no surprise, since the only boundary-parabolic $\SL(2,\C)$-representations of the figure-$8$ knot are abelian. To detect the geometric representation, we need to consider \emph{obstruction classes} (see Section~\ref{sec:ObstructionClasses} below).

\subsubsection{The figure-$8$ knot sister}
Consider the ideal triangulation of the figure-$8$ knot sister shown in Figure~\ref{fig:Fig8SisOri}. The Ptolemy variety $P_2(\T)$ is given by
\begin{equation}\label{eq:Fig8SisPtolemyUnsimplified}
\begin{aligned}[c]
c_{03,0}c_{12,0}+c_{01,0}c_{23,0}&=c_{02,0}c_{13,0},\\
c_{03,1}c_{12,1}+c_{01,1}c_{23,1}&=c_{02,1}c_{13,1},
\end{aligned}
\qquad 
\begin{aligned}[c]
c_{01,0}=-c_{03,0}=c_{23,0}=-c_{01,1}=c_{03,1}=-c_{23,1}\\
c_{02,0}=-c_{12,0}=c_{13,0}=-c_{02,1}=c_{12,1}=-c_{13,1}
\end{aligned}
\end{equation}
Selecting representatives $\varepsilon_{23,0}$ and $\varepsilon_{13,0}$ for the $2$ edge cycles, $P_2(\T)\in\C^2$ is given by
\begin{equation}\label{eq:Fig8SisPtolemyTrivial}
c_{23,0}c_{13,0}+c_{23,0}^2=c_{13,0}^2,\qquad c_{23,0}c_{13,0}+c_{23,0}^2=c_{13,0}^2.
\end{equation}
This is equivalent to
\begin{equation}
x^2-x-1=0,\qquad x=\frac{c_{13,0}}{c_{23,0}}.
\end{equation}
%By Theorem~\ref{thm:PtolemyFieldEqualsTraceField}, $P_2(\T)$ detects two representations with trace field $\Q(\sqrt{5})$.

%When fixing the decoration $c_{23,0}$ is set to $1$.

\begin{figure}[htpb]
\centering
\begin{minipage}[c]{0.48\textwidth}
\scalebox{0.57}{\begingroup
 \setlength{\unitlength}{0.8pt}
 \begin{picture}(522.27557,212.41415)
 \put(0,0){\includegraphics{figures_gen/Fig8Triangulation_included}}

\definecolor{inkcol1}{rgb}{0.0,0.0,0.0}
   \put(94.449251,12.219258){\rotatebox{360.0}{\makebox(0,0)[tl]{\strut{}{
    \begin{minipage}[h]{129.560208pt}
\textcolor{inkcol1}{\LARGE{$c_{02,0}$}}\\
\end{minipage}}}}}%


\definecolor{inkcol1}{rgb}{0.0,0.0,0.0}
   \put(33.020681,129.362118){\rotatebox{360.0}{\makebox(0,0)[tl]{\strut{}{
    \begin{minipage}[h]{129.560208pt}
\textcolor{inkcol1}{\LARGE{$c_{01,0}$}}\\
\end{minipage}}}}}%


\definecolor{inkcol1}{rgb}{0.0,0.0,0.0}
   \put(243.734971,116.504978){\rotatebox{360.0}{\makebox(0,0)[tl]{\strut{}{
    \begin{minipage}[h]{129.560208pt}
\textcolor{inkcol1}{\LARGE{$c_{23,0}$}}\\
\end{minipage}}}}}%


\definecolor{inkcol1}{rgb}{0.0,0.0,0.0}
   \put(172.306401,219.362128){\rotatebox{360.0}{\makebox(0,0)[tl]{\strut{}{
    \begin{minipage}[h]{129.560208pt}
\textcolor{inkcol1}{\LARGE{$c_{13,0}$}}\\
\end{minipage}}}}}%


\definecolor{inkcol1}{rgb}{0.0,0.0,0.0}
   \put(170.877821,118.647838){\rotatebox{360.0}{\makebox(0,0)[tl]{\strut{}{
    \begin{minipage}[h]{129.560208pt}
\textcolor{inkcol1}{\LARGE{$c_{12,0}$}}\\
\end{minipage}}}}}%


\definecolor{inkcol1}{rgb}{0.0,0.0,0.0}
   \put(88.734971,123.647838){\rotatebox{360.0}{\makebox(0,0)[tl]{\strut{}{
    \begin{minipage}[h]{129.560208pt}
\textcolor{inkcol1}{\LARGE{$c_{03,0}$}}\\
\end{minipage}}}}}%


\definecolor{inkcol1}{rgb}{0.0,0.0,0.0}
   \put(410.877831,220.076408){\rotatebox{360.0}{\makebox(0,0)[tl]{\strut{}{
    \begin{minipage}[h]{129.560208pt}
\textcolor{inkcol1}{\LARGE{$c_{13,1}$}}\\
\end{minipage}}}}}%


\definecolor{inkcol1}{rgb}{0.0,0.0,0.0}
   \put(488.020691,129.362118){\rotatebox{360.0}{\makebox(0,0)[tl]{\strut{}{
    \begin{minipage}[h]{129.560208pt}
\textcolor{inkcol1}{\LARGE{$c_{23,1}$}}\\
\end{minipage}}}}}%


\definecolor{inkcol1}{rgb}{0.0,0.0,0.0}
   \put(338.734971,12.219258){\rotatebox{360.0}{\makebox(0,0)[tl]{\strut{}{
    \begin{minipage}[h]{129.560208pt}
\textcolor{inkcol1}{\LARGE{$c_{02,1}$}}\\
\end{minipage}}}}}%


\definecolor{inkcol1}{rgb}{0.0,0.0,0.0}
   \put(282.306391,135.790688){\rotatebox{360.0}{\makebox(0,0)[tl]{\strut{}{
    \begin{minipage}[h]{129.560208pt}
\textcolor{inkcol1}{\LARGE{$c_{01,1}$}}\\
\end{minipage}}}}}%


\definecolor{inkcol1}{rgb}{0.0,0.0,0.0}
   \put(410.877821,118.647828){\rotatebox{360.0}{\makebox(0,0)[tl]{\strut{}{
    \begin{minipage}[h]{129.560208pt}
\textcolor{inkcol1}{\LARGE{$c_{12,1}$}}\\
\end{minipage}}}}}%


\definecolor{inkcol1}{rgb}{0.0,0.0,0.0}
   \put(330.877831,125.790678){\rotatebox{360.0}{\makebox(0,0)[tl]{\strut{}{
    \begin{minipage}[h]{129.560208pt}
\textcolor{inkcol1}{\LARGE{$c_{03,1}$}}\\
\end{minipage}}}}}%

 \end{picture}
\endgroup
}
\captionsetup{width=10cm}
\caption{Ordered triangulation of the figure $8$ knot. The signs indicate the non-trivial obstruction class.}\label{fig:Fig8Ordered}
\end{minipage}	
\hspace{2mm}
\begin{minipage}[c]{0.48\textwidth}
\scalebox{0.57}{\begingroup
 \setlength{\unitlength}{0.8pt}
 \begin{picture}(522.27557,212.41415)
 \put(0,0){\includegraphics{figures_gen/Fig8SisterTriangulation_included}}

\definecolor{inkcol1}{rgb}{0.0,0.0,0.0}
   \put(173.458011,219.841498){\rotatebox{360.0}{\makebox(0,0)[tl]{\strut{}{
    \begin{minipage}[h]{129.560208pt}
\textcolor{inkcol1}{\LARGE{$c_{13,0}$}}\\
\end{minipage}}}}}%


\definecolor{inkcol1}{rgb}{0.0,0.0,0.0}
   \put(418.925081,221.861818){\rotatebox{360.0}{\makebox(0,0)[tl]{\strut{}{
    \begin{minipage}[h]{129.560208pt}
\textcolor{inkcol1}{\LARGE{$c_{13,1}$}}\\
\end{minipage}}}}}%


\definecolor{inkcol1}{rgb}{0.0,0.0,0.0}
   \put(337.276041,10.739918){\rotatebox{360.0}{\makebox(0,0)[tl]{\strut{}{
    \begin{minipage}[h]{129.560208pt}
\textcolor{inkcol1}{\LARGE{$c_{02,1}$}}\\
\end{minipage}}}}}%


\definecolor{inkcol1}{rgb}{0.0,0.0,0.0}
   \put(488.625611,124.887168){\rotatebox{360.0}{\makebox(0,0)[tl]{\strut{}{
    \begin{minipage}[h]{129.560208pt}
\textcolor{inkcol1}{\LARGE{$c_{23,1}$}}\\
\end{minipage}}}}}%


\definecolor{inkcol1}{rgb}{0.0,0.0,0.0}
   \put(277.503721,134.988688){\rotatebox{360.0}{\makebox(0,0)[tl]{\strut{}{
    \begin{minipage}[h]{129.560208pt}
\textcolor{inkcol1}{\LARGE{$c_{01,1}$}}\\
\end{minipage}}}}}%


\definecolor{inkcol1}{rgb}{0.0,0.0,0.0}
   \put(413.874321,114.785638){\rotatebox{360.0}{\makebox(0,0)[tl]{\strut{}{
    \begin{minipage}[h]{129.560208pt}
\textcolor{inkcol1}{\LARGE{$c_{12,1}$}}\\
\end{minipage}}}}}%


\definecolor{inkcol1}{rgb}{0.0,0.0,0.0}
   \put(330.031651,125.897318){\rotatebox{360.0}{\makebox(0,0)[tl]{\strut{}{
    \begin{minipage}[h]{129.560208pt}
\textcolor{inkcol1}{\LARGE{$c_{03,1}$}}\\
\end{minipage}}}}}%


\definecolor{inkcol1}{rgb}{0.0,0.0,0.0}
   \put(92.645821,10.862478){\rotatebox{360.0}{\makebox(0,0)[tl]{\strut{}{
    \begin{minipage}[h]{129.560208pt}
\textcolor{inkcol1}{\LARGE{$c_{02,0}$}}\\
\end{minipage}}}}}%


\definecolor{inkcol1}{rgb}{0.0,0.0,0.0}
   \put(32.036651,130.948078){\rotatebox{360.0}{\makebox(0,0)[tl]{\strut{}{
    \begin{minipage}[h]{129.560208pt}
\textcolor{inkcol1}{\LARGE{$c_{01,0}$}}\\
\end{minipage}}}}}%


\definecolor{inkcol1}{rgb}{0.0,0.0,0.0}
   \put(239.117931,111.755188){\rotatebox{360.0}{\makebox(0,0)[tl]{\strut{}{
    \begin{minipage}[h]{129.560208pt}
\textcolor{inkcol1}{\LARGE{$c_{23,0}$}}\\
\end{minipage}}}}}%


\definecolor{inkcol1}{rgb}{0.0,0.0,0.0}
   \put(87.595051,123.877008){\rotatebox{360.0}{\makebox(0,0)[tl]{\strut{}{
    \begin{minipage}[h]{129.560208pt}
\textcolor{inkcol1}{\LARGE{$c_{03,0}$}}\\
\end{minipage}}}}}%


\definecolor{inkcol1}{rgb}{0.0,0.0,0.0}
   \put(173.335481,115.795808){\rotatebox{360.0}{\makebox(0,0)[tl]{\strut{}{
    \begin{minipage}[h]{129.560208pt}
\textcolor{inkcol1}{\LARGE{$c_{12,0}$}}\\
\end{minipage}}}}}%

 \end{picture}
\endgroup
}
\captionsetup{width=10cm}
\caption{Oriented triangulation of the figure $8$ knot sister. The signs indicate the non-trivial obstruction class.}\label{fig:Fig8SisOri}
\end{minipage}					
%\caption{The census triangulation of the figure 8 knot sister. No vertex ordering exists making all face pairings order preserving.}\label{PtolemyFigure}
\end{figure}
\begin{remark}
Note that for ordered triangulations, the identification relations~\eqref{eq:IdentificationRelations} do not involve minus signs. The triangulation in Figure~\ref{fig:Fig8Ordered} is not oriented.



\end{remark}
\subsection{Obstruction classes}\label{sec:ObstructionClasses}
Each class in $H^2(\widehat M;\Z/2\Z)$ can be represented by a $\Z/2\Z$-valued $2$-cocycle on $\widehat M$, i.e.~an assignment of a sign to each face of $\T$.  %There is a Ptolemy variety $P^{\sigma}(\T)$ for each class $\sigma\in H^2(\widehat M;\Z/2\Z)$. 
\begin{definition} Let $\sigma$ be a $\Z/2\Z$-valued $2$-cocycle on $\widehat M$. The \emph{Ptolemy variety} for $\sigma$ is defined as in Definition~\ref{def:PtolemyVariety}, but with the Ptolemy relation replaced by
\begin{equation}\label{eq:PtolemyRelationObstruction}
\sigma_{0,k}\sigma_{3,k}c_{03,k}c_{12,k}+\sigma_{0,k}\sigma_{1,k}c_{01,k}c_{23,k}=\sigma_{0,k}\sigma_{2,k}c_{02,k}c_{13,k},
\end{equation}
where $\sigma_{i,k}$ is the sign of the face of $\Delta_k$ opposite vertex $i$.
\end{definition}
\begin{remark} Multiplying $\sigma$ by a coboundary $\delta(\tau)$ corresponds to multiplying the Ptolemy coordinate of a one-cell $e$ by $\tau(e)$ (see~\cite{GaroufalidisThurstonZickert} for details). Hence, up to canonical isomorphism, the Ptolemy variety $P^\sigma_2(\T)$ only depends on the cohomology class of $\sigma$. The Ptolemy variety $P_2(\T)$ is the Ptolemy variety for the trivial obstruction class.
\end{remark}

\subsubsection{Examples}
In both examples above, $H^2(\widehat M;\Z/2\Z)=\Z/2\Z$, and the non-trivial obstruction class $\sigma$ is indicated in Figures~\ref{fig:Fig8Ordered} and~\ref{fig:Fig8SisOri}.

For the Figure-$8$ knot, $P^{\sigma}_2(\T)$ is given by
\begin{equation}\label{eq:Fig8PtolemyNontrivial}
-c_{23,0}c_{13,0}+c_{23,0}^2=-c_{13,0}^2,\qquad -c_{13,0}c_{23,0}+c_{13,0}^2=-c_{23,0}^2,
\end{equation}
which is equivalent to
\begin{equation}
x^2-x+1=0,\qquad x=\frac{c_{13,0}}{c_{23,0}}.
\end{equation} 
The corresponding representations are the geometric representation and its conjugate.

For the Figure 8 knot sister, the Ptolemy variety becomes
\begin{equation}\label{eq:Fig8SisPtolemyNontrivial}
-c_{23,0}c_{13,0}-c_{23,0}^2=c_{13,0}^2,\qquad -c_{23,0}c_{13,0}-c_{23,0}^2=c_{13,0}^2,
\end{equation}
which is equivalent to
\begin{equation}
x^2+x+1=0,\qquad x=\frac{c_{13,0}}{c_{23,0}}.
\end{equation} 


\section{The diagonal action}
Fix an ordering of the $1$-cells of $\T$ and of the cusps of $M$. As mentioned in Remark~\ref{rm:PickRepresentative} the Ptolemy variety can be regarded as a subset of the ambient space $\C^e$. 



Let $T=(\C^*)^v$ be a torus whose coordinates are indexed by the cusps of $M$. There is a natural action of $T$ on $P^{\sigma}_2(\T)$ defined as follows:
For $x=(x_1,\dots,x_v)\in T$ and $c=(c_1,\dots c_e)\in P^{\sigma}_2(\T)$ define a Ptolemy assignment $cx$ by
\begin{equation}
(xc)_i=x_jx_kc_i,
\end{equation}
where $j$ and $k$ (possibly $j=k$) are the cusps joined by the $i$th edge cycle. The action is thus determined entirely by the $1$-skeleton of $\widehat M$.

\begin{remark}\label{rm:IntrinsicDiagonalAction} There is a more intrinsic definition of this action in terms of decorations: 
Each vertex of $\Mtildehat$ determines a cusp of $M$, and if $D$ is a decoration taking a vertex $w$ to $gP$, the decoration $xD$ takes $w$ to $g\left(\begin{smallmatrix}x_i&0\\0&x_i^{-1}\end{smallmatrix}\right)P$ where $i$ is the cusp determined by $w$. The fact that the two definitions agree under the one-one correspondence \eqref{eq:1to1Obstruction} is an immediate consequence of the relationship given in Figure~\ref{fig:DecToPtolemy}.
\end{remark}

\subsection{The reduced Ptolemy varieties}
\begin{definition}
The \emph{reduced Ptolemy variety} $P^\sigma_2(\T)_\rd$ is the quotient $P^\sigma_2(\T)\big/T$. 
\end{definition}
Let $\mathcal O$ be the coordinate ring of $P^\sigma_2(\T)$, and let $\mathcal O^T$ be the ring of invariants. By geometric invariant theory, the reduced Ptolemy variety is a variety whose coordinate ring is isomorphic to $\mathcal O^T$.

For $i=0,1$, let $C_i$ denote the free abelian group generated by the \emph{unoriented} $i$-cells of $\widehat M$, and consider the maps (first studied by Neumann~\cite{NeumannComb})
\begin{equation}\label{eq:AlphaAlphaStar}
\alpha\colon C_0\to C_1, \qquad \alpha^*\colon C_1\to C_0,
\end{equation}
where $\alpha$ takes a $0$-cell to a sum of its incident $1$-cells, and $\alpha^*$ takes a $1$-cell to the sum of its end points. The maps $\alpha$ and $\alpha^*$ are dual under the canonical identifications $C_i\cong C_i^*$. Also, $\alpha$ is injective, and $\alpha^*$ has cokernel of order $2$ (see \cite{NeumannComb}.)

%First studied by Neumann~\cite{NeumannComb}, and generalized to $n>2$ in Garoufalidis--Zickert~\cite{GaroufalidisZickert}. The cokernel of $\alpha^*$ is $\Z/2\Z$.

The following is an elementary consequence of the definition of the diagonal action.
\begin{lemma}
The diagonal action $P_2^\sigma(\T)$, and the induced action on the coordinate ring $\mathcal O$ of $P_2^\sigma(\T)$ are given, respectively, by
\begin{equation}\label{eq:TActionOnPtolemys}
(xc)_i=(\prod_{j=1}^v x_j^{\alpha_{ij}})c_i, \qquad
x(c^w)=\prod_{j=1}^v x_j^{\alpha^*(w)_j}c^w,
\end{equation}
where $c^w$ is the monomial $c_1^{w_1}\cdots c_e^{w_e}\in\mathcal O$, $w\in\Z^e$.\qed
\end{lemma}

\begin{corollary}\label{cor:RingOfInvs}
Let $w_1,\dots,w_{e-v}$ be a basis for $\Ker(\alpha^*)$. The monomials
$c^{w_1},\dots,c^{w_{e-v}}$ generate $\mathcal O^T$.
\qed
\end{corollary}

\subsubsection{Examples}
Suppose the $1$-skeleton of $\widehat M$ looks like in Figure~\ref{fig:Before} (this is in fact the $1$-skeleton of the census triangulation of the Whitehead link complement.) We have
\begin{equation}
\alpha^*=\begin{pmatrix}2&1&1&0\\0&1&1&2\end{pmatrix}
\end{equation}
and the action of $(x_1,x_2)$ on a Ptolemy assignment $c$ is given in Figure~\ref{fig:After}.
\begin{figure}[htpb]
\centering
\begin{minipage}[c]{0.48\textwidth}
\scalebox{0.85}{\begingroup
 \setlength{\unitlength}{0.8pt}
 \begin{picture}(343.69867,70.627403)
 \put(0,0){\includegraphics{figures_gen/PtolemyWhBefore_included}}

\definecolor{inkcol1}{rgb}{0.0,0.0,0.0}
   \put(118.053087,74.427329){\rotatebox{360.0}{\makebox(0,0)[tl]{\strut{}{
    \begin{minipage}[h]{127.655896pt}
\textcolor{inkcol1}{\Large{$c_2$}}\\
\end{minipage}}}}}%


\definecolor{inkcol1}{rgb}{0.0,0.0,0.0}
   \put(-1.041283,34.559439){\rotatebox{360.0}{\makebox(0,0)[tl]{\strut{}{
    \begin{minipage}[h]{127.655896pt}
\textcolor{inkcol1}{\Large{$c_1$}}\\
\end{minipage}}}}}%


\definecolor{inkcol1}{rgb}{0.0,0.0,0.0}
   \put(229.672997,36.702289){\rotatebox{360.0}{\makebox(0,0)[tl]{\strut{}{
    \begin{minipage}[h]{127.655896pt}
\textcolor{inkcol1}{\Large{$c_4$}}\\
\end{minipage}}}}}%


\definecolor{inkcol1}{rgb}{0.0,0.0,0.0}
   \put(116.101587,1.702289){\rotatebox{360.0}{\makebox(0,0)[tl]{\strut{}{
    \begin{minipage}[h]{127.655896pt}
\textcolor{inkcol1}{\Large{$c_3$}}\\
\end{minipage}}}}}%


\definecolor{inkcol1}{rgb}{0.0,0.0,0.0}
   \put(76.037957,18.737429){\rotatebox{360.0}{\makebox(0,0)[tl]{\strut{}{
    \begin{minipage}[h]{127.655896pt}
\textcolor{inkcol1}{\Large{$1$}}\\
\end{minipage}}}}}%


\definecolor{inkcol1}{rgb}{0.0,0.0,0.0}
   \put(158.180807,19.451729){\rotatebox{360.0}{\makebox(0,0)[tl]{\strut{}{
    \begin{minipage}[h]{127.655896pt}
\textcolor{inkcol1}{\Large{$2$}}\\
\end{minipage}}}}}%

 \end{picture}
\endgroup
}
\captionsetup{width=10cm}
\caption{Ptolemy assignment.}\label{fig:Before}
\end{minipage}	
\begin{minipage}[c]{0.48\textwidth}
\scalebox{0.85}{\begingroup
 \setlength{\unitlength}{0.8pt}
 \begin{picture}(350.29688,73.169228)
 \put(0,0){\includegraphics{figures_gen/PtolemyWhAfter_included}}

\definecolor{inkcol1}{rgb}{0.0,0.0,0.0}
   \put(124.481647,76.747044){\rotatebox{360.0}{\makebox(0,0)[tl]{\strut{}{
    \begin{minipage}[h]{127.655896pt}
\textcolor{inkcol1}{\Large{$x_1x_2c_2$}}\\
\end{minipage}}}}}%


\definecolor{inkcol1}{rgb}{0.0,0.0,0.0}
   \put(-1.041283,38.307734){\rotatebox{360.0}{\makebox(0,0)[tl]{\strut{}{
    \begin{minipage}[h]{127.655896pt}
\textcolor{inkcol1}{\Large{$x_1^2c_1$}}\\
\end{minipage}}}}}%


\definecolor{inkcol1}{rgb}{0.0,0.0,0.0}
   \put(248.958707,41.879164){\rotatebox{360.0}{\makebox(0,0)[tl]{\strut{}{
    \begin{minipage}[h]{127.655896pt}
\textcolor{inkcol1}{\Large{$x_2^2c_4$}}\\
\end{minipage}}}}}%


\definecolor{inkcol1}{rgb}{0.0,0.0,0.0}
   \put(124.673007,2.593434){\rotatebox{360.0}{\makebox(0,0)[tl]{\strut{}{
    \begin{minipage}[h]{127.655896pt}
\textcolor{inkcol1}{\Large{$x_1x_2c_3$}}\\
\end{minipage}}}}}%


\definecolor{inkcol1}{rgb}{0.0,0.0,0.0}
   \put(97.466537,22.485734){\rotatebox{360.0}{\makebox(0,0)[tl]{\strut{}{
    \begin{minipage}[h]{127.655896pt}
\textcolor{inkcol1}{\Large{$1$}}\\
\end{minipage}}}}}%


\definecolor{inkcol1}{rgb}{0.0,0.0,0.0}
   \put(174.609377,23.200024){\rotatebox{360.0}{\makebox(0,0)[tl]{\strut{}{
    \begin{minipage}[h]{127.655896pt}
\textcolor{inkcol1}{\Large{$2$}}\\
\end{minipage}}}}}%

 \end{picture}
\endgroup
}
\captionsetup{width=10cm}
\caption{The diagonal action of $(x_1,x_2)$.}\label{fig:After}
\end{minipage}					
%\caption{The census triangulation of the figure 8 knot sister. No vertex ordering exists making all face pairings order preserving.}\label{PtolemyFigure}
\end{figure}
The kernel of $\alpha^*$ is generated by $(0,-2,0,1)^t$ and $(-1,1,1,0)^t$, so we have
\begin{equation}
\mathcal O^T=\langle c_2^{-2}c_4,c_1^{-1}c_2c_3\rangle.
\end{equation}
Also note that in each of the examples in Section~\ref{sec:PtolemyVariety}, $x\in\mathcal O^T$.

For computations we need a more explicit description of the reduced Ptolemy variety.
\begin{definition}
Let $T\colon \Z^n\to\Z^m$ be a homomorphism. We say that $T$ is \emph{basic} if there exists a subset $J$ of $\{e_1,\dots,e_n\}$ such that $T$ maps $\Span(J)$ isomorphically onto the image of $T$. Elements of such a set $J$ are called \emph{basic generators} for $T$.
%We call such generators of $\Imag(T)$ \emph{basic generators}.
\end{definition}
 We identify $C_1$ and $C_0$ with $\Z^e$ and $\Z^v$, respectively.
%Rationally, every homomorphism is basic, but not integrally.
\begin{proposition}\label{prop:Basic}
The map $\alpha^*\colon C_1\to C_0$ is basic.\qed
\end{proposition}
The proof will be relegated to Section~\ref{sec:AlphaStarBasic}, where we shall also give explicit basic generators.
\begin{proposition}\label{prop:ReducedPtolemyVariety}
Let $\varepsilon_{i_1},\dots,\varepsilon_{i_v}$ be basic generators for $\alpha^*$. 
The ring of invariants $\mathcal O^T$ is isomorphic to $\C[c_1,\dots,c_e]$ modulo the Ptolemy relations and the relations $c_{i_1}=\cdots=c_{i_v}=1$, i.e.~the reduced Ptolemy variety is isomorphic to the subset of $P_2^\sigma(\T)$ where the Ptolemy coordinates of the basic generators are $1$.
\end{proposition}
\begin{proof}
Let $w_1,\dots,w_{e-v}$ be a basis for $\Ker(\alpha^*)$. Hence, $w_1,\dots,w_{e-v}$ and $\varepsilon_{i_1},\dots,\varepsilon_{i_v}$ generate $C_1$. We can thus uniquely express each $c_i$ as a monomial in the $w_j$'s and the $c_{i_j}$'s. The result now follows from Corollary~\ref{cor:RingOfInvs}.
\end{proof}

\begin{remark}
This is how the Ptolemy varieties are computed in SnapPy.
\end{remark}
%\subsection{Gluing equations}
%Say something about gluing equations. Mention that for each set $I$ of basic generators, the gluing equations are generated by the gluing equations of edges not in $I$.
\subsection{Shapes and gluing equations}
One can assign to each simplex a \emph{shape}
\begin{equation}
z=\sigma_3\sigma_2\frac{c_{03}c_{12}}{c_{02}c_{13}}\in\C\setminus\{0,1\},
\end{equation}
and one can show (see \cite{GaroufalidisThurstonZickert,GaroufalidisGoernerZickert}) that these satisfy Thurston's gluing equations.
For the geometric representation of a cusped hyperbolic manifold, the shape field (field generated by the shapes) is equal to the invariant trace field, which is in general smaller than the trace field, see Maclachlan--Reid~\cite{MaclachlanReid}.
\begin{remark}
Note that the shapes are elements in $\mathcal O^T$.
\end{remark}


\subsection{Proof that $\alpha^*$ is basic}\label{sec:AlphaStarBasic}
Since $\alpha^*$ has cokernel of order $2$, it is enough to prove that there is a set of columns of $\alpha^*$ forming a matrix with determinant $\pm 2$.
Recall that the columns of $\alpha^*$ corresponds to $1$-cells of $\T$. We shall thus consider graphs in the $1$-skeleton of $\widehat M$. We recall some basic results from graph theory. All graphs are assumed to be connected.
%We split the proof into three cases. 
\begin{definition}
The \emph{incidence matrix} of a graph $G$ with vertices $v_1,\dots,v_k$ and edges $\varepsilon_1,\dots,\varepsilon_l$ is the $k\times l$ matrix $I_G$
whose $ij$th entry is $1$ if $v_i$ is incident to $\varepsilon_j$, and $0$ otherwise.
\end{definition}

\begin{lemma}\label{lemma:IncidenceMatrix}
The rank of $I_G$ is $k-1$. If $G$ is a tree, $I_G$ is a $k\times (k-1)$ matrix, and removing any row gives a matrix with determinant $\pm 1$.\qed
\end{lemma}
%\begin{proof}
%Clearly, this holds for a tree with $2$ vertices and $1$ edge. Let $\widetilde I_T$ be $I_T$ with a row removed. If $v_i$ is a vertex whose only incident edge is $\varepsilon_j$, removing row $i$ and column $j$ from $\widetilde I_T$ does not change the determinant. The result now follows by induction.
%\end{proof}

\subsubsection{Case $1$: A single cusp}
In this case the result is trivial. The matrix representation for $\alpha^*$ is $\begin{pmatrix}2&\cdots&2\end{pmatrix}$.

\subsubsection{Case $2$: Multiple cusps, self edges}
Suppose $\widehat M$ has a self edge $\varepsilon_1$ (an edge joining a cusp to itself), and consider the graph $G$ consisting of the union of $\varepsilon_1$ with a maximal tree $T$ (see Figure~\ref{fig:TreeWithCycle}). The columns of $\alpha^*$ corresponding to the edges of $G$ then form the matrix 
\begin{equation}
B=\left(\begin{array}{c|c}
2& \\\cline{1-1}
 0 &\raisebox{10pt}{{$I_T$}}
\end{array}\right)
\end{equation}

which, by Lemma~\ref{lemma:IncidenceMatrix}, has determinant $\pm 2$.
\subsubsection{Case $3$: Multiple cusps, no self edges}
Pick a face with edges $\varepsilon_1,\varepsilon_2,\varepsilon_3$, and add edges to form a graph $G$ such that $G\setminus \varepsilon_1$ is a maximal tree (see Figure~\ref{fig:TreeWithThreeCycle}). The corresponding columns form the matrix

\begin{equation}
C=I_G=\left(\begin{array}{c|c}
1&\\ 
0&I_T\\
1&\\\cline{1-1} 
 0&
\end{array}\right)
\end{equation}
which, by Lemma~\ref{lemma:IncidenceMatrix}, has determinant $\pm 2$.
This concludes the proof that $\alpha^*$ is basic.

Note that 
\begin{equation}
\det(B)=\det\begin{pmatrix}2&1\\&1\end{pmatrix}=2,\qquad \det(C)=\det\begin{pmatrix}1&1&\\0&1&1\\1&&1\end{pmatrix}=2,
\end{equation}
i.e.~only the edges and vertices shown in Figures~\ref{fig:TreeWithCycle} and \ref{fig:TreeWithThreeCycle} contribute to the determinant.
\begin{figure}[htb]
\centering
\begin{minipage}[c]{0.45\textwidth}
\centering
\scalebox{1}{\begingroup
 \setlength{\unitlength}{0.8pt}
 \begin{picture}(230.76326,61.12545)
 \put(0,0){\includegraphics{figures_gen/TreeWithCycle_included}}

\definecolor{inkcol1}{rgb}{0.0,0.0,0.0}
   \put(10.47901,38.692798){\rotatebox{360.0}{\makebox(0,0)[tl]{\strut{}{
    \begin{minipage}[h]{129.560208pt}
\textcolor{inkcol1}{\large{$\varepsilon_1$}}\\
\end{minipage}}}}}%


\definecolor{inkcol1}{rgb}{0.0,0.0,0.0}
   \put(89.224612,47.580128){\rotatebox{360.0}{\makebox(0,0)[tl]{\strut{}{
    \begin{minipage}[h]{129.560208pt}
\textcolor{inkcol1}{\large{$\varepsilon_2$}}\\
\end{minipage}}}}}%


\definecolor{inkcol1}{rgb}{0.0,0.0,0.0}
   \put(69.944022,24.332028){\rotatebox{360.0}{\makebox(0,0)[tl]{\strut{}{
    \begin{minipage}[h]{129.560208pt}
\textcolor{inkcol1}{\large{$1$}}\\
\end{minipage}}}}}%


\definecolor{inkcol1}{rgb}{0.0,0.0,0.0}
   \put(118.944022,24.082028){\rotatebox{360.0}{\makebox(0,0)[tl]{\strut{}{
    \begin{minipage}[h]{129.560208pt}
\textcolor{inkcol1}{\large{$2$}}\\
\end{minipage}}}}}%

 \end{picture}
\endgroup
}
\end{minipage}
\hfill
\begin{minipage}[c]{0.45\textwidth}
\centering
\begingroup
 \setlength{\unitlength}{0.8pt}
 \begin{picture}(190.10619,94.455963)
 \put(0,0){\includegraphics{figures_gen/TreeWithThreeCycle_included}}

\definecolor{inkcol1}{rgb}{0.0,0.0,0.0}
   \put(10.518997,55.14957){\rotatebox{360.0}{\makebox(0,0)[tl]{\strut{}{
    \begin{minipage}[h]{111.71428pt}
\textcolor{inkcol1}{\large{$\varepsilon_2$}}\\
\end{minipage}}}}}%


\definecolor{inkcol1}{rgb}{0.0,0.0,0.0}
   \put(38.390127,17.1655){\rotatebox{360.0}{\makebox(0,0)[tl]{\strut{}{
    \begin{minipage}[h]{111.71428pt}
\textcolor{inkcol1}{\large{$\varepsilon_3$}}\\
\end{minipage}}}}}%


\definecolor{inkcol1}{rgb}{0.0,0.0,0.0}
   \put(61.919517,54.75215){\rotatebox{360.0}{\makebox(0,0)[tl]{\strut{}{
    \begin{minipage}[h]{111.71428pt}
\textcolor{inkcol1}{\large{$\varepsilon_1$}}\\
\end{minipage}}}}}%


\definecolor{inkcol1}{rgb}{0.0,0.0,0.0}
   \put(48.523387,90.53842){\rotatebox{360.0}{\makebox(0,0)[tl]{\strut{}{
    \begin{minipage}[h]{111.71428pt}
\textcolor{inkcol1}{\large{$1$}}\\
\end{minipage}}}}}%


\definecolor{inkcol1}{rgb}{0.0,0.0,0.0}
   \put(17.053287,18.26919){\rotatebox{360.0}{\makebox(0,0)[tl]{\strut{}{
    \begin{minipage}[h]{111.71428pt}
\textcolor{inkcol1}{\large{$2$}}\\
\end{minipage}}}}}%


\definecolor{inkcol1}{rgb}{0.0,0.0,0.0}
   \put(58.031337,18.59325){\rotatebox{360.0}{\makebox(0,0)[tl]{\strut{}{
    \begin{minipage}[h]{111.71428pt}
\textcolor{inkcol1}{\large{$3$}}\\
\end{minipage}}}}}%

 \end{picture}
\endgroup

\end{minipage}
\\
\begin{minipage}[t]{0.45\textwidth}
\caption{Tree $G$ with $1$-cycle; $G\setminus\varepsilon_1$ is a maximal tree.}\label{fig:TreeWithCycle}
\end{minipage}
\begin{minipage}[t]{0.45\textwidth}
\caption{Tree $G$ with $3$-cycle; $G\setminus\varepsilon_1$ is a maximal tree.}\label{fig:TreeWithThreeCycle}
\end{minipage}
\end{figure}

\begin{remark}
Trees with $1$- or $3$-cycles are also used in~\cite[Sec.4.6]{GHRS} to study index structures.
\end{remark}


\section{The Ptolemy field and the trace field}
\subsection{Explicit description of the Ptolemy field}
By Proposition~\ref{prop:ReducedPtolemyVariety} any $c\in P^\sigma_2(\T)$ is equivalent to a Ptolemy assignment $c^\prime$ whose coordinates for a set of basic generators $\varepsilon_{i_1},\dots,\varepsilon_{i_v}$ is $1$. In particular, it follows that the Ptolemy field (see Definition~\ref{def:PtolemyField}) of $c\in P^\sigma_2(\T)$ is given by
\begin{equation}\label{eq:kc}
k_c=k_{c^{\prime}}=\Q\big(\big\{c_{\varepsilon_1},\dots,c_{\varepsilon_e}\big\}\big).
\end{equation}





\begin{definition}
Let $\rho:\pi_1(M)\to\PSL(2,\C)$ be a representation. The \emph{trace field} of $\rho$ is the field generated by the traces of elements in the image. 
We denote it $k_\rho$.
\end{definition}
Our main result is the following. We defer the proof to Section~\ref{sec:ProofOfMainThm}.
\begin{theorem}\label{thm:PtolemyEqTrace} Let $c\in P^\sigma_2(\T)_\rd$. If the corresponding generic representation $\rho$ of $\pi_1(M)$ in $\PSL(2,\C)$ is boundary-non-trivial, the Ptolemy field of $c$ equals the trace field of $\rho$.\qed
\end{theorem}
\begin{remark}
Note that if $c\in P^\sigma_2(\T)_\rd$ is in a degree $0$ component, the Ptolemy field is a number field. %The Ptolemy fields for all zero-dimensional components are listed in unhyperbolic.
\end{remark}
%Note that the trace field is a subfield of the Ptolemy field. This follows immediately from Figure~\ref{}, which shows that the representation is defined over the Ptolemy field.
\subsection{The setup of the proof}
Since the natural cocycle is given in terms of the Ptolemy coordinates, it follows that $\rho$ is defined over the Ptolemy field. Hence, the trace field is a subfield of the Ptolemy field.

Fix a maximal tree $G$ with $1$ or $3$-cycle as in Figures~\ref{fig:TreeWithCycle} or \ref{fig:TreeWithThreeCycle}. As explained in Section~\ref{sec:AlphaStarBasic} the edges of $G$ are basic generators of $\alpha^*$. We may thus assume without loss of generality that the Ptolemy coordinates $c_i$ of the edges $\varepsilon_i$ of $G$ are $1$. By~\eqref{eq:kc}, it is thus enough to show that the Ptolemy coordinates of the remaining $1$-cells are in the trace field.


Let $\gamma$ denote the (lifted) natural cocycle of $c$. Then $\gamma$ assigns to each edge path $p$ in $\widehat M$ a matrix $\gamma(p)\in\SL(2,\C)$. Let
\begin{equation}\label{eq:AlphaBeta}
\alpha(a)=\begin{pmatrix}0&-a^{-1}\\a&0\end{pmatrix},\qquad \beta(b)=\begin{pmatrix}1&b\\0&1\end{pmatrix}.
\end{equation}
As shown in Figure~\ref{fig:PtolemyToCocycle}, $\gamma$ takes long and short edges to elements of the form $\alpha(a)$ and $\beta(b)$, respectively, where $a$ and $b$ are given in terms of the Ptolemy coordinates.

Since $\rho$ is boundary-non trivial, there exists for each cusp $i$ of $M$ a peripheral loop $M_i$ with $\gamma(M_i)\in P$ non-trivial. We shall here refer to such loops as \emph{non-trivial}. Fix such non-trivial loops $M_i$, once and for all, and let $m_i\neq 0$ be such that $\gamma(M_i)=\beta(m_i)$.
For any edge path $p$ with end point on a cusp $i$ we can alter $M_i$ by a conjugation if necessary (this does not change $m_i$) so that $p$ is composable with $M_i$.

\subsection{Proof for one cusp}
We first prove Theorem~\ref{thm:PtolemyEqTrace} in the case where there is only one cusp. In this case, all edges are self edges, and $T$ consists of a single edge $\varepsilon_1$.
\begin{lemma}\label{lemma:BootStrapEdgeOneCusp} For any self edge $\varepsilon$, we have $m_1c_\varepsilon\in k_\rho$. 
\end{lemma}
\begin{proof}
Let $X_1$ be a peripheral path such that $X_1\varepsilon$ is a loop (see Figure~\ref{fig:SelfEdgeCase}), and let $x_1$ be such that $\gamma(X_1)=\beta(x_1)$. We have
\begin{equation}\label{eq:OneCuspTraceRel}
\Tr(\gamma(X_1\varepsilon))=\Tr(\beta(x_1)\alpha(c_\varepsilon)) =\Tr\left(\begin{pmatrix}1&x_1\\0&1\end{pmatrix}\begin{pmatrix}0&-c_\varepsilon^{-1}\\c_\varepsilon&0\end{pmatrix}\right)=x_1c_\varepsilon\in k_\rho.
\end{equation}
Applying the same computation to the loop $X_1 M_1 \varepsilon$ yields
\begin{equation}
\Tr(\beta(x_1)\beta(m_1)\alpha(c_\varepsilon))=(x_1+m_1)c_\varepsilon\in k_\rho,
\end{equation}
and the result follows.
\end{proof}
Since the Ptolemy coordinate of $\varepsilon_1$ is $1$, it follows that $m_1\in k_\rho$. Since all edges are self edges, we have $c_\varepsilon\in k_\rho$ for all $1$-cells $\varepsilon$. This concludes the proof in the one cusped case.


\begin{figure}[htb]
\begin{center}
\hspace{-10mm}
\begin{minipage}[c]{0.3\textwidth}
\scalebox{0.8}{\begingroup
 \setlength{\unitlength}{0.8pt}
 \begin{picture}(235.52596,245.00639)
 \put(0,0){\includegraphics{figures_gen/SelfEdgeBootStrap_included}}

\definecolor{inkcol1}{rgb}{0.0,0.0,0.0}
   \put(157.61284,157.186559){\rotatebox{360.0}{\makebox(0,0)[tl]{\strut{}{
    \begin{minipage}[h]{98.516592pt}
\textcolor{inkcol1}{\large{$\beta(m_1)$}}\\
\end{minipage}}}}}%


\definecolor{inkcol1}{rgb}{0.0,0.0,0.0}
   \put(74.40285,147.405567){\rotatebox{360.0}{\makebox(0,0)[tl]{\strut{}{
    \begin{minipage}[h]{69.262756pt}
\textcolor{inkcol1}{\large{$X_1$}}\\
\end{minipage}}}}}%


\definecolor{inkcol1}{rgb}{0.0,0.0,0.0}
   \put(75.28559,124.596669){\rotatebox{360.0}{\makebox(0,0)[tl]{\strut{}{
    \begin{minipage}[h]{98.516592pt}
\textcolor{inkcol1}{\large{$\beta(x_1)$}}\\
\end{minipage}}}}}%


\definecolor{inkcol1}{rgb}{0.0,0.0,0.0}
   \put(152.05261,76.736547){\rotatebox{360.0}{\makebox(0,0)[tl]{\strut{}{
    \begin{minipage}[h]{95.93536pt}
\textcolor{inkcol1}{\large{$\varepsilon$}}\\
\end{minipage}}}}}%


\definecolor{inkcol1}{rgb}{0.0,0.0,0.0}
   \put(177.02256,66.75248){\rotatebox{360.0}{\makebox(0,0)[tl]{\strut{}{
    \begin{minipage}[h]{154.012832pt}
\textcolor{inkcol1}{\large{$\alpha(c_\varepsilon)$}}\\
\end{minipage}}}}}%


\definecolor{inkcol1}{rgb}{0.0,0.0,0.0}
   \put(126.74744,147.557807){\rotatebox{360.0}{\makebox(0,0)[tl]{\strut{}{
    \begin{minipage}[h]{69.262756pt}
\textcolor{inkcol1}{\large{$M_1$}}\\
\end{minipage}}}}}%

 \end{picture}
\endgroup
}
\end{minipage}
\hspace{2mm}
\begin{minipage}[c]{0.6\textwidth}
\scalebox{0.8}{\begingroup
 \setlength{\unitlength}{0.8pt}
 \begin{picture}(447.41534,169.52409)
 \put(0,0){\includegraphics{figures_gen/TwoCuspEdge_included}}

\definecolor{inkcol1}{rgb}{0.0,0.0,0.0}
   \put(192.092684,23.022627){\rotatebox{360.0}{\makebox(0,0)[tl]{\strut{}{
    \begin{minipage}[h]{176.040232pt}
\textcolor{inkcol1}{\large{$\alpha(c_{\varepsilon})$}}\\
\end{minipage}}}}}%


\definecolor{inkcol1}{rgb}{0.0,0.0,0.0}
   \put(191.10446,72.425037){\rotatebox{360.0}{\makebox(0,0)[tl]{\strut{}{
    \begin{minipage}[h]{176.040232pt}
\textcolor{inkcol1}{\large{$\alpha(-c_{\varepsilon})$}}\\
\end{minipage}}}}}%


\definecolor{inkcol1}{rgb}{0.0,0.0,0.0}
   \put(321.108055,71.183027){\rotatebox{360.0}{\makebox(0,0)[tl]{\strut{}{
    \begin{minipage}[h]{176.040232pt}
\textcolor{inkcol1}{\large{$\beta(-m_{j})$}}\\
\end{minipage}}}}}%


\definecolor{inkcol1}{rgb}{0.0,0.0,0.0}
   \put(76.758924,70.086057){\rotatebox{360.0}{\makebox(0,0)[tl]{\strut{}{
    \begin{minipage}[h]{176.040232pt}
\textcolor{inkcol1}{\large{$\beta(m_i)$}}\\
\end{minipage}}}}}%

 \end{picture}
\endgroup
}
\end{minipage}
\hfill
\\
\begin{minipage}[t]{0.32\textwidth}
\captionsetup{width=1\textwidth}
\caption{Self edge.}\label{fig:SelfEdgeCase}
\end{minipage}
\hspace{2mm}
\begin{minipage}[t]{0.64\textwidth}
\captionsetup{width=\textwidth}
\caption[Bla]{Edge between cusps.}\label{fig:TwoCuspCase}
\end{minipage}
\end{center}
\end{figure}

\subsection{The general case} \label{sec:ProofOfMainThm}The general case follows the same strategy, but is more complicated since it involves edge paths between multiple cusps. %The key step is the result below, which is easily proved by induction. We shall not need the explicit formula.
%Assume that $T$ is a maximal tree with $1$-cycle (see~\eqref{}). The proof for the case when $T$ is a maximal tree with $3$-cycle is virtually identical and left to the reader.

\begin{lemma}\label{lemma:BootStrapEdge}
If $\varepsilon$ is a self edge from cusp $i$ to itself, $m_ic_\varepsilon\in k_\rho$
\end{lemma}
\begin{proof}
The proof is identical to that of Lemma~\ref{lemma:BootStrapEdgeOneCusp}.
\end{proof}

\begin{lemma}\label{lemma:connectTwoCusp}
If two (distinct) cusps $i$ and $j$ are joined by an edge $\varepsilon$ in $G$, we have
\begin{equation}m_i m_j\in k_\rho.
\end{equation} 
\end{lemma}

\begin{proof}
Consider the loop $\varepsilon_j\bar{M}_j\bar{\varepsilon_j}M_i$ shown in Figure~\ref{fig:TwoCuspCase}. A simple computation shows that
\begin{equation}
\Tr(\alpha(c_\varepsilon)\beta(-m_{j})\alpha(-c_{\varepsilon})\beta(m_i))=2+m_i m_{j} c_{\varepsilon}^2.
\end{equation}
Since $\varepsilon\in T$, $c_\varepsilon=1$, and the result follows.
\end{proof}
More generally, the following holds.
\begin{lemma}\label{lemma:AllMiInKRho}
We have $m_i\in k_\rho$ for all cusps $i$.
\end{lemma}

\begin{proof}
If $G$ is a tree with 1-cycle, then $c_1=1$, so Lemma~\ref{lemma:BootStrapEdge} implies that $m_1\in k_\rho$. Inductively applying Lemma~\ref{lemma:connectTwoCusp} for the edge $\varepsilon_j$ connecting cusp $i=j-1$ and $j$ implies the result. If $G$ is a tree with $3$-cycle, the Ptolemy coordinates $c_1,c_2$ and $c_3$ are $1$, so the edges of the face are labeled by $\alpha(1)$ and $\beta(-1)$ only (see Figure~\ref{fig:PtolemyToCocycle}). Inserting the peripheral loops $M_i$ as in Figure~\ref{fig:TreeWithThreeCycleCase}, we obtain
\begin{equation}
\Tr\big(\beta(-1)\beta(m_1)\alpha(1)\beta(-1)\beta(m_2)\alpha(1)\beta(-1)\beta(m_3)\alpha(1)\big)\in k_\rho
\end{equation}
By an elementary computation, the trace equals
\begin{equation}
m_1m_2m_3-m_1m_2-m_2m_3-m_3m_1+2\in k_\rho.
\end{equation}
By Lemma~\ref{lemma:connectTwoCusp}, $m_im_j\in k_\rho$, so $m_1\in k_\rho$. The result now follows as above by inductively applying Lemma~\ref{lemma:connectTwoCusp}.
\end{proof}

Let $\varepsilon$ be an arbitrary $1$-cell. If $\varepsilon$ is a self edge, Lemmas~\ref{lemma:BootStrapEdge} and \ref{lemma:AllMiInKRho} imply that $c_\varepsilon\in k_\rho$. Otherwise, there exists an edge path $p$ in the maximal tree $G\setminus\varepsilon_1$, such that $p*\varepsilon$ is a loop in $\widehat M$. By relabeling the cusps and edges if necessary, we may assume that $p=\varepsilon_{i+1}\!*\!\varepsilon_{i+2}\!*\cdots\!*\!\varepsilon_{j}$, where $\varepsilon_{k}$ goes from cusp $k-1$ to cusp $k$. Pick peripheral paths $X_{k}$ on cusp $k$ connecting the ends (in $M$, not $\widehat M$) of edges $\varepsilon_{k}$ and $\varepsilon_{k+1}$ (see Figure~\ref{fig:ArbitraryEdgeCase}). We obtain a loop that can be composed with arbitrary powers of the peripheral loops $M_{i}, \dots, M_{j}$. We thus obtain the following traces (where $b_k\in\Z$)
\begin{equation}
\Tr(\beta(x_{i}+b_{i}m_{i})\alpha(c_{i+1})\beta(x_{i+1}+b_{i+1}m_{i+1})\alpha(c_{i+2})\cdots \beta(x_{j}+b_jm_j)\alpha(c_\varepsilon))\in k_\rho.
\end{equation}

\begin{figure}[htb]
\begin{center}
\hspace{-10mm}
\begin{minipage}[c]{0.4\textwidth}
\scalebox{0.8}{\begingroup
 \setlength{\unitlength}{0.8pt}
 \begin{picture}(340.52692,277.40344)
 \put(0,0){\includegraphics{figures_gen/ThreeCuspCycle_included}}

\definecolor{inkcol1}{rgb}{0.0,0.0,0.0}
   \put(182.5451029,184.4897843){\rotatebox{360.0}{\makebox(0,0)[tl]{\strut{}{
    \begin{minipage}[h]{365.37732pt}
\textcolor{inkcol1}{\large{$\alpha(1)$}}\\
\end{minipage}}}}}%


\definecolor{inkcol1}{rgb}{0.0,0.0,0.0}
   \put(178.759021,57.27859){\rotatebox{360.0}{\makebox(0,0)[tl]{\strut{}{
    \begin{minipage}[h]{365.37732pt}
\textcolor{inkcol1}{\large{$\alpha(1)$}}\\
\end{minipage}}}}}%


\definecolor{inkcol1}{rgb}{0.0,0.0,0.0}
   \put(44.74235,127.858837){\rotatebox{360.0}{\makebox(0,0)[tl]{\strut{}{
    \begin{minipage}[h]{365.37732pt}
\textcolor{inkcol1}{\large{$\alpha(1)$}}\\
\end{minipage}}}}}%


\definecolor{inkcol1}{rgb}{0.0,0.0,0.0}
   \put(64.56978,217.033079){\rotatebox{360.0}{\makebox(0,0)[tl]{\strut{}{
    \begin{minipage}[h]{365.37732pt}
\textcolor{inkcol1}{\large{$\beta(-1)$}}\\
\end{minipage}}}}}%


\definecolor{inkcol1}{rgb}{0.0,0.0,0.0}
   \put(247.718131,123.025426){\rotatebox{360.0}{\makebox(0,0)[tl]{\strut{}{
    \begin{minipage}[h]{365.37732pt}
\textcolor{inkcol1}{\large{$\beta(-1)$}}\\
\end{minipage}}}}}%


\definecolor{inkcol1}{rgb}{0.0,0.0,0.0}
   \put(70.71563,22.56969){\rotatebox{360.0}{\makebox(0,0)[tl]{\strut{}{
    \begin{minipage}[h]{365.37732pt}
\textcolor{inkcol1}{\large{$\beta(-1)$}}\\
\end{minipage}}}}}%


\definecolor{inkcol1}{rgb}{0.0,0.0,0.0}
   \put(173.530509,267.189497){\rotatebox{360.0}{\makebox(0,0)[tl]{\strut{}{
    \begin{minipage}[h]{365.37732pt}
\textcolor{inkcol1}{\large{$\beta(m_1)$}}\\
\end{minipage}}}}}%


\definecolor{inkcol1}{rgb}{0.0,0.0,0.0}
   \put(279.517293,29.66383){\rotatebox{360.0}{\makebox(0,0)[tl]{\strut{}{
    \begin{minipage}[h]{365.37732pt}
\textcolor{inkcol1}{\large{$\beta(m_2)$}}\\
\end{minipage}}}}}%


\definecolor{inkcol1}{rgb}{0.0,0.0,0.0}
   \put(16.27348,89.252617){\rotatebox{360.0}{\makebox(0,0)[tl]{\strut{}{
    \begin{minipage}[h]{365.37732pt}
\textcolor{inkcol1}{\large{$\beta(m_3)$}}\\
\end{minipage}}}}}%


\definecolor{inkcol1}{rgb}{0.0,0.0,0.0}
   \put(99.697986,48.33525){\rotatebox{360.0}{\makebox(0,0)[tl]{\strut{}{
    \begin{minipage}[h]{365.37732pt}
\textcolor{inkcol1}{\large{$X_3$}}\\
\end{minipage}}}}}%


\definecolor{inkcol1}{rgb}{0.0,0.0,0.0}
   \put(99.372859,188.800326){\rotatebox{360.0}{\makebox(0,0)[tl]{\strut{}{
    \begin{minipage}[h]{365.37732pt}
\textcolor{inkcol1}{\large{$X_1$}}\\
\end{minipage}}}}}%


\definecolor{inkcol1}{rgb}{0.0,0.0,0.0}
   \put(215.163868,122.470312){\rotatebox{360.0}{\makebox(0,0)[tl]{\strut{}{
    \begin{minipage}[h]{365.37732pt}
\textcolor{inkcol1}{\large{$X_2$}}\\
\end{minipage}}}}}%


\definecolor{inkcol1}{rgb}{0.0,0.0,0.0}
   \put(92.135175,125.210053){\rotatebox{360.0}{\makebox(0,0)[tl]{\strut{}{
    \begin{minipage}[h]{365.37732pt}
\textcolor{inkcol1}{\large{$\varepsilon_1$}}\\
\end{minipage}}}}}%


\definecolor{inkcol1}{rgb}{0.0,0.0,0.0}
   \put(168.905141,160.137307){\rotatebox{360.0}{\makebox(0,0)[tl]{\strut{}{
    \begin{minipage}[h]{365.37732pt}
\textcolor{inkcol1}{\large{$\varepsilon_2$}}\\
\end{minipage}}}}}%


\definecolor{inkcol1}{rgb}{0.0,0.0,0.0}
   \put(162.520277,74.00083){\rotatebox{360.0}{\makebox(0,0)[tl]{\strut{}{
    \begin{minipage}[h]{365.37732pt}
\textcolor{inkcol1}{\large{$\varepsilon_3$}}\\
\end{minipage}}}}}%


\definecolor{inkcol1}{rgb}{0.0,0.0,0.0}
   \put(23.21089,59.00358){\rotatebox{360.0}{\makebox(0,0)[tl]{\strut{}{
    \begin{minipage}[h]{365.37732pt}
\textcolor{inkcol1}{\large{$M_3$}}\\
\end{minipage}}}}}%


\definecolor{inkcol1}{rgb}{0.0,0.0,0.0}
   \put(138.154585,247.301417){\rotatebox{360.0}{\makebox(0,0)[tl]{\strut{}{
    \begin{minipage}[h]{365.37732pt}
\textcolor{inkcol1}{\large{$M_1$}}\\
\end{minipage}}}}}%


\definecolor{inkcol1}{rgb}{0.0,0.0,0.0}
   \put(250.314966,60.01249){\rotatebox{360.0}{\makebox(0,0)[tl]{\strut{}{
    \begin{minipage}[h]{365.37732pt}
\textcolor{inkcol1}{\large{$M_2$}}\\
\end{minipage}}}}}%

 \end{picture}
\endgroup
}
\end{minipage}
\begin{minipage}[c]{0.6\textwidth}
\scalebox{0.8}{\begingroup
 \setlength{\unitlength}{0.8pt}
 \begin{picture}(467.052,214.4512)
 \put(0,0){\includegraphics{figures_gen/ManyCuspCase_included}}

\definecolor{inkcol1}{rgb}{0.0,0.0,0.0}
   \put(82.821724,84.7552665){\rotatebox{360.0}{\makebox(0,0)[tl]{\strut{}{
    \begin{minipage}[h]{365.37732pt}
\textcolor{inkcol1}{\large{$\alpha(c_{i+1})$}}\\
\end{minipage}}}}}%


\definecolor{inkcol1}{rgb}{0.0,0.0,0.0}
   \put(13.5915,138.34793){\rotatebox{360.0}{\makebox(0,0)[tl]{\strut{}{
    \begin{minipage}[h]{365.37732pt}
\textcolor{inkcol1}{\large{$\beta(x_i)$}}\\
\end{minipage}}}}}%


\definecolor{inkcol1}{rgb}{0.0,0.0,0.0}
   \put(80.164288,185.658862){\rotatebox{360.0}{\makebox(0,0)[tl]{\strut{}{
    \begin{minipage}[h]{365.37732pt}
\textcolor{inkcol1}{\large{$\beta(m_i)$}}\\
\end{minipage}}}}}%


\definecolor{inkcol1}{rgb}{0.0,0.0,0.0}
   \put(39.987221,109.856966){\rotatebox{360.0}{\makebox(0,0)[tl]{\strut{}{
    \begin{minipage}[h]{365.37732pt}
\textcolor{inkcol1}{\large{$X_i$}}\\
\end{minipage}}}}}%


\definecolor{inkcol1}{rgb}{0.0,0.0,0.0}
   \put(83.759839,106.586428){\rotatebox{360.0}{\makebox(0,0)[tl]{\strut{}{
    \begin{minipage}[h]{365.37732pt}
\textcolor{inkcol1}{\large{$\varepsilon_{i+1}$}}\\
\end{minipage}}}}}%


\definecolor{inkcol1}{rgb}{0.0,0.0,0.0}
   \put(52.483056,161.747852){\rotatebox{360.0}{\makebox(0,0)[tl]{\strut{}{
    \begin{minipage}[h]{365.37732pt}
\textcolor{inkcol1}{\large{$M_i$}}\\
\end{minipage}}}}}%


\definecolor{inkcol1}{rgb}{0.0,0.0,0.0}
   \put(175.90683,53.904203){\rotatebox{360.0}{\makebox(0,0)[tl]{\strut{}{
    \begin{minipage}[h]{365.37732pt}
\textcolor{inkcol1}{\large{$\alpha(c_\varepsilon)$}}\\
\end{minipage}}}}}%


\definecolor{inkcol1}{rgb}{0.0,0.0,0.0}
   \put(183.6829,28.394938){\rotatebox{360.0}{\makebox(0,0)[tl]{\strut{}{
    \begin{minipage}[h]{365.37732pt}
\textcolor{inkcol1}{\large{$\varepsilon$}}\\
\end{minipage}}}}}%


\definecolor{inkcol1}{rgb}{0.0,0.0,0.0}
   \put(181.757894,84.7552665){\rotatebox{360.0}{\makebox(0,0)[tl]{\strut{}{
    \begin{minipage}[h]{365.37732pt}
\textcolor{inkcol1}{\large{$\alpha(c_{i+2})$}}\\
\end{minipage}}}}}%


\definecolor{inkcol1}{rgb}{0.0,0.0,0.0}
   \put(99.761713,138.34793){\rotatebox{360.0}{\makebox(0,0)[tl]{\strut{}{
    \begin{minipage}[h]{365.37732pt}
\textcolor{inkcol1}{\large{$\beta(x_{i+1})$}}\\
\end{minipage}}}}}%


\definecolor{inkcol1}{rgb}{0.0,0.0,0.0}
   \put(179.100458,185.658862){\rotatebox{360.0}{\makebox(0,0)[tl]{\strut{}{
    \begin{minipage}[h]{365.37732pt}
\textcolor{inkcol1}{\large{$\beta(m_{i+1})$}}\\
\end{minipage}}}}}%


\definecolor{inkcol1}{rgb}{0.0,0.0,0.0}
   \put(132.5404125,109.856966){\rotatebox{360.0}{\makebox(0,0)[tl]{\strut{}{
    \begin{minipage}[h]{365.37732pt}
\textcolor{inkcol1}{\large{$X_{i+1}$}}\\
\end{minipage}}}}}%


\definecolor{inkcol1}{rgb}{0.0,0.0,0.0}
   \put(185.278645,106.586428){\rotatebox{360.0}{\makebox(0,0)[tl]{\strut{}{
    \begin{minipage}[h]{365.37732pt}
\textcolor{inkcol1}{\large{$\varepsilon_{i+2}$}}\\
\end{minipage}}}}}%


\definecolor{inkcol1}{rgb}{0.0,0.0,0.0}
   \put(153.546886,161.747852){\rotatebox{360.0}{\makebox(0,0)[tl]{\strut{}{
    \begin{minipage}[h]{365.37732pt}
\textcolor{inkcol1}{\large{$M_{i+1}$}}\\
\end{minipage}}}}}%


\definecolor{inkcol1}{rgb}{0.0,0.0,0.0}
   \put(300.82554,138.34793){\rotatebox{360.0}{\makebox(0,0)[tl]{\strut{}{
    \begin{minipage}[h]{365.37732pt}
\textcolor{inkcol1}{\large{$\beta(x_j)$}}\\
\end{minipage}}}}}%


\definecolor{inkcol1}{rgb}{0.0,0.0,0.0}
   \put(367.39833,185.658862){\rotatebox{360.0}{\makebox(0,0)[tl]{\strut{}{
    \begin{minipage}[h]{365.37732pt}
\textcolor{inkcol1}{\large{$\beta(m_j)$}}\\
\end{minipage}}}}}%


\definecolor{inkcol1}{rgb}{0.0,0.0,0.0}
   \put(325.0936,109.856966){\rotatebox{360.0}{\makebox(0,0)[tl]{\strut{}{
    \begin{minipage}[h]{365.37732pt}
\textcolor{inkcol1}{\large{$X_j$}}\\
\end{minipage}}}}}%


\definecolor{inkcol1}{rgb}{0.0,0.0,0.0}
   \put(341.31284,161.747852){\rotatebox{360.0}{\makebox(0,0)[tl]{\strut{}{
    \begin{minipage}[h]{365.37732pt}
\textcolor{inkcol1}{\large{$M_j$}}\\
\end{minipage}}}}}%


\definecolor{inkcol1}{rgb}{0.0,0.0,0.0}
   \put(362.85956,70.495264){\rotatebox{360.0}{\makebox(0,0)[tl]{\strut{}{
    \begin{minipage}[h]{365.37732pt}
\textcolor{inkcol1}{\large{Cusp $j$}}\\
\end{minipage}}}}}%


\definecolor{inkcol1}{rgb}{0.0,0.0,0.0}
   \put(277.1942,110.459941){\rotatebox{360.0}{\makebox(0,0)[tl]{\strut{}{
    \begin{minipage}[h]{365.37732pt}
\textcolor{inkcol1}{\large{$\varepsilon_{j}$}}\\
\end{minipage}}}}}%


\definecolor{inkcol1}{rgb}{0.0,0.0,0.0}
   \put(269.09307,84.5765278){\rotatebox{360.0}{\makebox(0,0)[tl]{\strut{}{
    \begin{minipage}[h]{365.37732pt}
\textcolor{inkcol1}{\large{$\alpha(c_{j})$}}\\
\end{minipage}}}}}%

 \end{picture}
\endgroup
}
\end{minipage}
\hfill
\\
\begin{minipage}[t]{0.35\textwidth}
\captionsetup{width=1\textwidth}
\caption{3-cycle case.}\label{fig:TreeWithThreeCycleCase}
\end{minipage}
\hspace{2mm}
\begin{minipage}[t]{0.5\textwidth}
\captionsetup{width=\textwidth}
\caption[Bla]{Arbitrary edge $\varepsilon$.}\label{fig:ArbitraryEdgeCase}
\end{minipage}
\end{center}
\end{figure}

It will be convenient to regard $\Tr(\beta(x_i)\alpha(c_{i+1})\beta(x_{i+1})\alpha(c_{i+2})\cdots \beta(x_j)\alpha(c_\varepsilon))$ as a function of variables $x_i$ (disregarding that the $x_i$'s are fixed expressions of the Ptolemy coordinates).
\begin{definition}
Given a function $f(x_1,\dots,x_r)$, let $\Delta_{i}f$ be the function given by
\begin{equation}
\Delta_{i}f(h)=f(x_1,\dots,x_i+h,\dots,x_r)-f(x_1,\dots,x_i,\dots,x_r).
\end{equation}
\end{definition}
The following is elementary.%; see~\eqref{eq:DeltaExample} for an example.
\begin{lemma}\label{lemma:Deltas} If $f(x_1,\dots,x_r)$ is a polynomial where the exponents of all variables $x_i$ are $0$ or $1$, and where the highest degree term is $ax_1x_2\cdots x_r$, we have
\begin{equation}\label{eq:Deltas}
\Delta_r\big(\cdots\Delta_2\big(\Delta_1f(h_1)\big)(h_2)\cdots\big)=ah_1h_2\dots h_r,
\end{equation}
and is thus independent of the $x_i$'s.\qed
\end{lemma}
If, e.g.,~$f(x_1,x_2)=x_1x_2$, we have
\begin{equation}\label{eq:DeltaExample}
\Delta_1f(h_1)=(x_1+h_1)x_2-x_1x_2=h_1x_2,\quad \Delta_2\big(\Delta_1f(h_1)\big)(h_2)=h_1(x_2+h_2)-h_1x_2=h_1h_2.
\end{equation}


\begin{lemma}\label{lemma:TraceFormula} Let $x_1,\dots,x_r$ be variables and $y_1,\dots,y_r$ be constants.
The expression
\begin{equation}
\Tr\big(\beta(x_1)\alpha(y_1)\cdots\beta(x_r)\alpha(y_r)\big)
\end{equation}
is a polynomial in the $x_i$'s whose unique highest degree term is $\prod_{i=1}^ry_i\prod_{i=1}^rx_i$. Moreover, for each monomial term, the exponent of each variable is either $1$ or $0$.
\end{lemma}
\begin{proof}
This follows by induction on $r$.
\end{proof}

Applying Lemmas \ref{lemma:TraceFormula} and \ref{lemma:Deltas} to the function
\begin{equation}
f(x_i,\dots,x_j)=\Tr\big(\beta(x_i)\alpha(c_{i+1})\beta(x_{i+1})\alpha(c_{i+2})\cdots \beta(x_j)\alpha(c_\varepsilon)\big),
\end{equation} 
we obtain 
\begin{equation}
(m_im_{i+1}\cdots m_jc_ic_{i+1}\cdots c_j)c_\varepsilon\in k_\rho.
\end{equation}
Since all $m_i$'s are in $k_\rho$ by Lemma~\ref{lemma:AllMiInKRho}, and all $c_i$'s are $1$ (since $\varepsilon_i\in T$), it follows that $c_\varepsilon$ is in $k_\rho$. This concludes the proof. 


\subsection{Proof of Proposition~\ref{prop:GenericBoundaryNonTrivial}}\label{sec:PtolemyToRepProof}
%This follows trivially from results in Garoufalidis--Thurston--Zickert~\cite{GaroufalidisThurstonZickert}.
The fact that $\mathcal R$ factors, follows from the fact that the diagonal action only changes the decoration (by diagonal elements, c.f.~Remark~\ref{rm:IntrinsicDiagonalAction}), not the representation.
Since the preimage of the right map in~\eqref{eq:1to1Obstruction} is parametrized by choices of lifts, i.e.~elements in $Z^1(\widehat M;\Z/2\Z)$, all that remains is to show that the only freedom in the choice of decoration of a boundary-non-trivial representation is the diagonal action. 
This follows from results in ~\cite{GaroufalidisThurstonZickert}: A decoration is an equivariant map
\begin{equation}
D\colon \widehat{\widetilde{M}}^{(0)}\to \PSL(2,\C)\big/P,
\end{equation}
and is thus determined by its image of lifts $\widetilde e_1,\dots,\widetilde e_v$ of the cusps of $M$. The freedom in the choice of $D(\widetilde e_i)$ is the choice of a coset $gP$ satisfying that $g\rho(\Stab(\widetilde e_i))g^{-1}\subset P$, where $\Stab(\widetilde e_i)\subset\pi_1(M)$ is the stabilizer of $\widetilde e_i$, i.e.~a peripheral subgroup corresponding to cusp $i$. Hence, if $\rho(\Stab(\widetilde e_i))$ is non-trivial, the freedom is right multiplication by a diagonal matrix (if it is trivial, any coset works). Hence, if $\rho$ is boundary-non-trivial, the only freedom in choosing a decoration is the diagonal action. 
%For $n>2$, there are a lot more possibilities. 


%\section{Computations}
%\label{sec.computations}

\section{Ptolemy varieties for $n>2$}\label{sec:HigherPtolemy}
Many of our results generalize in a straightforward way to the higher rank Ptolemy varieties $P_n(\T)$. We recall the definition of these below, and refer to~\cite{GaroufalidisThurstonZickert,GaroufalidisGoernerZickert} for details.

We identify all simplices of $\T$ with a standard simplex
\begin{equation}
\Delta^3_n = \left\{(x_0,x_1,x_2,x_3) \in \R^4\bigm\vert 0\leq x_i \leq n,\enspace x_0 + x_1 + x_2 + x_3 = n\right\}
\end{equation}
and regard $\widehat M$ as a quotient of a disjoint union $\coprod_{k=1}^s\Delta^3_{n,k}$, with a copy $\Delta^3_{n,k}$ of $\Delta^3_n$ for each simplex $k$ of $\T$. Define $\Delta^3_n(\Z)=\Delta^3_n\cap\Z^4$ and $\dot\Delta^3_n(\Z)$ to be $\Delta^3_n(\Z)$ with the four vertex points removed.
A point in $M$ in the image of $\coprod_{k=1}^s\dot\Delta^3_{n,k}(\Z)$ is called an \emph{integral point} of $M$.

\subsection{Definition of the Ptolemy variety}
Assign to each $(t,k)\in\Delta^3_{n,k}(\Z)$ a \emph{Ptolemy coordinate} $c_{t,k}$.
For each simplex $k$, we have $\vert\Delta_{n-2}(\Z)\vert=\binom{n+1}{3}$ \emph{Ptolemy relations}
\begin{equation}
c_{\alpha+1001,k}c_{\alpha+0110,k}+c_{\alpha+1100,k}c_{\alpha+0011,k}=c_{\alpha+1010,k}c_{\alpha+0101,k},
\qquad \alpha\in\Delta_{n-2}(\Z)
\end{equation}
as well as \emph{identification relations},
\begin{equation}\label{eq:IdentificationRelationsN}
c_{t,k}=\pm c_{t^{\prime},k^{\prime}},\qquad\text{when}\qquad (t,k)\sim (t^{\prime},k^{\prime}).
\end{equation}
%for each integral point of $M$.
\begin{remark}The signs in~\ref{eq:IdentificationRelationsN} depend in a non-trivial way on the face pairings (see~\cite{GaroufalidisGoernerZickert}). For ordered triangulations the signs are always positive. As in Remark~\ref{rm:PickRepresentative} we can eliminate the identification relations by selecting a representative of each integral point of $M$.
\end{remark}
\begin{definition} The \emph{Ptolemy variety} $P_n(\T)$ is the subset of the affine algebraic set defined by the Ptolemy- and identification relations, consisting of the points where all Ptolemy coordinates are non-zero.
\end{definition}
For general $n$, we denote the group of upper triangular matrices with $1$ on the diagonal by $N$ (instead of $P$). As in~\ref{eq:summaryof1to1} we have %(see~\cite{GaroufalidisThurstonZickert} for a proof) 
\begin{equation}\label{eq:summaryof1to1GeneralN}
\xymatrix{\left\{\txt{Points in\\$P_n(\T)$}\right\}\ar@{<->}[r]^-{1-1}&\left\{\txt{Natural $(\SL(n,\C),N)$-\\cocycles on $M$}\right\}\ar@{<->}[r]^-{1-1}&\left\{\txt{Generically decorated\\
$(\SL(n,\C),N)$-representations}\right\}.}
\end{equation}

\subsection{The diagonal action}
Let $D$ be the group of diagonal matrices in $\SL(n,\C)$. We identify $D$ with the torus $(\C^*)^{n-1}$ via the identification
\begin{equation}
(\C^*)^{n-1}\to D,\qquad (a_1,\dots,a_{n-1})\mapsto \diag(a_1,a_2/a_1,\dots,a_{n-1}/a_{n-2})
\end{equation}
As in Remark~\ref{rm:IntrinsicDiagonalAction}, we have a diagonal action of the torus $T=D^{v}$ on the set of decorated representations, where $(D_1,\dots,D_v)\in T$ acts by replacing the coset $gN$ assigned to a vertex $w$ by $gD_iN$, where $i$ is the cusp corresponding to $w$. The corresponding action on $P_n(\T)$ is described in Lemma~\ref{lemma:ActionOnPtolemy} below.

Let $C_1^n$ be the group generated by the integral points of $M$, and let $C^n_0=C_0\otimes \Z^{n-1}$. In Garoufalidis--Zickert~\cite{GaroufalidisZickert} we defined maps
\begin{equation}
\alpha\colon C^n_0\to C^n_1,\qquad \alpha^*\colon C^n_1\to C^n_0
\end{equation}
generalizing~\eqref{eq:AlphaAlphaStar}. The map $\alpha^*$ takes an integral point $(t,k)$ to $\sum x_i\otimes e_{t_i}$, where $x_i$ is the cusp determined by vertex $i$ of simplex $k$. We shall not need the definition of $\alpha$.
\begin{lemma}[Garoufalidis--Zickert~\cite{GaroufalidisZickert}]\label{lemma:CokernelN}
The map $\alpha^*$ is surjective with cokernel $\Z/n\Z$.
\end{lemma}
By selecting an ordering of the natural generators of $C^n_0$ and $C^n_1$, we regard $\alpha$ and $\alpha^*$ as matrices. The following is an elementary consequence of~\eqref{eq:summaryof1to1GeneralN}.
\begin{lemma}\label{lemma:ActionOnPtolemy}
The diagonal action of $T=(\C^*)^{v(n-1)}$ on $P_n(\T)$ and the corresponding action on the coordinate ring $\mathcal O$ of $P_n(\T)$ are given, respectively, by 
\begin{equation}
\pushQED{\qed}
(xc)_t=(\prod_{j=1}^{v(n-1)} x_j^{\alpha_{ij}})c_t,\qquad x(c^w)=\prod_{j=1}^{v(n-1)} x_j^{\alpha^*(w)_j}c^w.\qedhere
\popQED
%x(c^w)=\prod_{j=1}^v x_j^{\alpha^*(w)_i}c^w,
\end{equation}
\end{lemma}
\begin{corollary} The ring of invariants $\mathcal O^T$ is generated by $c^{w_1},\dots,c^{w_r}$, where $w_1,\dots,w_r$ are a basis for $\Ker(\alpha^*)$, and $r=\rank(C^n_1)-\rank(C^n_0)$.
\end{corollary}
\begin{definition}
The \emph{Ptolemy field} of a Ptolemy assignment $c\in P_n(\T)$ is defined as
\begin{equation}
k_c=\Q(c^{w_1},\dots,c^{w_r})
\end{equation}
where $w_1,\dots w_r$ are (integral) generators of $\Ker(\alpha^*)$. 
\end{definition}
The following is proved in Section~\ref{sec:AlphaStarBasicGeneral}.
\begin{proposition}
The map $\alpha^*\colon C^n_1\to C^n_0$ is basic.\qed
\end{proposition}
\begin{corollary}
Let $p_1,\dots,p_{(n-1)v}$ be integral points that are basic generators of $C^n_1$. The ring $\mathcal O^T$ is generated by the Ptolemy relations together with the relations $c_{p_1}=\cdots=c_{p_{(n-1)v}}=1$. Equivalently, the reduced Ptolemy variety is isomorphic to the subvariety of $P_n(\T)$ consisting of Ptolemy assignments with $c_{p_i}=1$.
\end{corollary}
\begin{proof}
This follows the proof of Proposition~\ref{prop:ReducedPtolemyVariety} word by word.
\end{proof}
\begin{remark}
This is how the Ptolemy varieties and Ptolemy fields at~\cite{Unhyperbolic} are computed.
\end{remark}
\subsection{Representations}
\begin{definition}\label{def:BoundaryNonDeg}
Let $\rho$ be an $(\SL(n,\C),N)$-representation, and let $I_i$ denote the image of the peripheral subgroup corresponding to cusp $i$. We say that $\rho$ is \emph{boundary-non-degenerate} if each $I_i$ has an element whose Jordan canonical form has a single (maximal) Jordan block.
\end{definition}
\begin{proposition}
The map 
\begin{equation}
\mathcal R\colon P_n(\T)_\rd\to \big\{(\SL(n,\C),N)\text{-representations}\big\}
\big/\Conj.
\end{equation}
maps onto the generic representations, and the preimage of a generic boundary-non-degenerate representation consists of a single point.
\end{proposition}
\begin{proof}
The proof is identical to the proof in Section~\ref{sec:PtolemyToRepProof} for $n=2$.
\end{proof}
\begin{conjecture}
The Ptolemy field of a generic, boundary-non-degenerate representation is equal to its trace field.
\end{conjecture}

\begin{remark}Much of the theory also works for $\PSL(n,\C)$-representations by means of obstruction classes in $H^2(\widehat M;\Z/n\Z)$. When $n$ is even, obstruction classes in $H^2(\widehat M;\Z/2\Z)$ were defined in~\cite{GaroufalidisThurstonZickert} for representations in $p\SL(n,\C)=\SL(n,\C)/\pm I$. 
For $\PSL(n,\C)$-representations, both the Ptolemy field and the trace field are only defined up to $n$th roots of unity.
The generalized obstruction classes are used on the website~\cite{Unhyperbolic} and will be explained in a forthcoming publication. 
%Goerner--Zickert~\cite{GoernerZickert}
\end{remark}

\subsection{Proof that $\alpha^*$ is basic}\label{sec:AlphaStarBasicGeneral}
By Lemma~\ref{lemma:CokernelN} we need to prove the existence of integral points such that the corresponding columns of $\alpha^*$ form a matrix of determinant $\pm n$. As in Section~\ref{sec:AlphaStarBasic} we split the proof into three cases.

\subsubsection{Basic matrix algebra}
Let $I_k$ be the identity matrix, $R_{k}$ the sparse matrix whose first row contains entirely of $1$'s, $S_{k}$ the sparse matrix whose lower diagonal consists of $1$'s ($S_1=0$), and $T_k$ the sparse matrix whose lower right entry is $1$. The index $k$ denotes that the matrices are $k\times k$.
\begin{equation}
R_3=\begin{pmatrix}1&1&1\\&&\\&&\end{pmatrix},\quad S_3=\begin{pmatrix}&&\\1&\\&1\end{pmatrix},\quad T_3=\begin{pmatrix}&&\\&&\\&&1\end{pmatrix}
\end{equation}
%We shall need the following elementary results:
\begin{lemma}\label{lemma:DeterminantIdentity} We have
\begin{equation}
\det(I_k+R_k-S_k)=k+1,\qquad \det(I_k+R_k+T_k-S_k)=2k+1.
\end{equation}
\end{lemma}
\begin{proof}
This follows e.g.~by expanding the determinant using the last column. The matrices $I_k+R_k-S_k$ are shown below for $k=1,2,3$ and $4$.
\begin{equation}
\begin{pmatrix}2\end{pmatrix},\qquad \begin{pmatrix}2&1\\-1&1\end{pmatrix},\qquad \begin{pmatrix}2&1&1\\-1&1&\\&-1&1\end{pmatrix},\qquad\begin{pmatrix}2&1&1&1\\-1&1&&\\&-1&1&\\&&-1&1\end{pmatrix}
\end{equation}
For $I_k+R_k+T_k-S_k$, the only difference is that the lower right entry is now $2$.
\end{proof}
\begin{lemma}\label{lemma:BlockDeterminant}
Let $A$, $B$, $C$, $D$ be $k\times k$, $k\times l$, $l\times k$, and $l\times l$ matrices, respectively, and let $M=\left(\begin{smallmatrix}A &B\\C&D\end{smallmatrix}\right)$. If $D$ is invertible, we have
\begin{equation}
\pushQED{\qed}
\det(M)=\det(D)\det(A-BD^{-1}C).\qedhere\popQED
\end{equation}
\end{lemma}
\begin{proof}
This follows from the identity $\left(\begin{smallmatrix}A &B\\C&D\end{smallmatrix}\right)=\left(\begin{smallmatrix}I &B\\0&D\end{smallmatrix}\right)\left(\begin{smallmatrix}A-BD^{-1}C &0\\D^{-1}C&I\end{smallmatrix}\right)$.
\end{proof}
  
\subsubsection{One cusp}
%Since there is only one cusp $x$, the generators of $C_0$ are $x\otimes e_i$, $i=1,\dots,n-1$.
Pick any face of $\T$ and consider the integral points shown in Figures~\ref{fig:OneCuspBasicOdd} and \ref{fig:OneCuspBasicEven}. Let $A_n$ be the $(n-1)\times(n-1)$ matrix formed by the corresponding columns of $\alpha^*$. The columns are ordered as shown in the figures, and the rows, i.e. the generators $x\otimes e_i$ of $C^n_0$, are ordered in the natural way (increasing in $i$). The following is an immediate consequence of the definition of $\alpha^*$.
\begin{lemma}
The matrix $A_n$ is given by
\begin{equation}
A_{2k+1}=\begin{pmatrix}I_k+R_k+T_k&I_k\\S_k&I_k\end{pmatrix},\qquad A_{2k}=\begin{pmatrix}2&0\cdots01&0\\0&I_{k-1}+R_{k-1}&I_{k-1}\\0&S_{k-1}&I_{k-1}\end{pmatrix}
\end{equation}
\end{lemma}
\begin{corollary}
The determinant of $A_n$ is $\pm n$.
\end{corollary}
\begin{proof}
This follows from Lemma~\ref{lemma:BlockDeterminant} and Lemma~\ref{lemma:DeterminantIdentity}.
\end{proof}

\begin{figure}[htpb]
\centering
\begin{minipage}[c]{0.48\textwidth}
\scalebox{0.8}{\begingroup
 \setlength{\unitlength}{0.8pt}
 \begin{picture}(268.6568,184.63654)
 \put(0,0){\includegraphics{figures_gen/OneCuspBasicOdd_included}}

\definecolor{inkcol1}{rgb}{0.0,0.0,0.0}
   \put(-5.561203,15.525296){\rotatebox{360.0}{\makebox(0,0)[tl]{\strut{}{
    \begin{minipage}[h]{127.655896pt}
\textcolor{inkcol1}{\Large{$x$}}\\
\end{minipage}}}}}%


\definecolor{inkcol1}{rgb}{0.0,0.0,0.0}
   \put(93.695257,184.072556){\rotatebox{360.0}{\makebox(0,0)[tl]{\strut{}{
    \begin{minipage}[h]{127.655896pt}
\textcolor{inkcol1}{\Large{$x$}}\\
\end{minipage}}}}}%


\definecolor{inkcol1}{rgb}{0.0,0.0,0.0}
   \put(154.044197,96.256546){\rotatebox{360.0}{\makebox(0,0)[tl]{\strut{}{
    \begin{minipage}[h]{127.655896pt}
\textcolor{inkcol1}{\large{$2k$}}\\
\end{minipage}}}}}%


\definecolor{inkcol1}{rgb}{0.0,0.0,0.0}
   \put(189.822007,16.244256){\rotatebox{360.0}{\makebox(0,0)[tl]{\strut{}{
    \begin{minipage}[h]{127.655896pt}
\textcolor{inkcol1}{\Large{$x$}}\\
\end{minipage}}}}}%


\definecolor{inkcol1}{rgb}{0.0,0.0,0.0}
   \put(183.001667,44.856476){\rotatebox{360.0}{\makebox(0,0)[tl]{\strut{}{
    \begin{minipage}[h]{127.655896pt}
\textcolor{inkcol1}{\large{$k\!+\!1$}}\\
\end{minipage}}}}}%


\definecolor{inkcol1}{rgb}{0.0,0.0,0.0}
   \put(173.498027,60.573316){\rotatebox{360.0}{\makebox(0,0)[tl]{\strut{}{
    \begin{minipage}[h]{127.655896pt}
\textcolor{inkcol1}{\large{$k\!+\!2$}}\\
\end{minipage}}}}}%


\definecolor{inkcol1}{rgb}{0.0,0.0,0.0}
   \put(113.036297,79.637116){\rotatebox{360.0}{\makebox(0,0)[tl]{\strut{}{
    \begin{minipage}[h]{127.655896pt}
\textcolor{inkcol1}{\large{$k$}}\\
\end{minipage}}}}}%


\definecolor{inkcol1}{rgb}{0.0,0.0,0.0}
   \put(93.121647,117.219426){\rotatebox{360.0}{\makebox(0,0)[tl]{\strut{}{
    \begin{minipage}[h]{127.655896pt}
\textcolor{inkcol1}{\large{$2$}}\\
\end{minipage}}}}}%


\definecolor{inkcol1}{rgb}{0.0,0.0,0.0}
   \put(85.782287,132.176086){\rotatebox{360.0}{\makebox(0,0)[tl]{\strut{}{
    \begin{minipage}[h]{127.655896pt}
\textcolor{inkcol1}{\large{$1$}}\\
\end{minipage}}}}}%

 \end{picture}
\endgroup
}
\captionsetup{width=10cm}
\caption{Basic generators, $n=2k+1$.}\label{fig:OneCuspBasicOdd}
\end{minipage}		
\begin{minipage}[c]{0.48\textwidth}
\scalebox{0.8}{\begingroup
 \setlength{\unitlength}{0.8pt}
 \begin{picture}(268.6568,184.63655)
 \put(0,0){\includegraphics{figures_gen/OneCuspBasicEven_included}}

\definecolor{inkcol1}{rgb}{0.0,0.0,0.0}
   \put(99.549357,135.585706){\rotatebox{360.0}{\makebox(0,0)[tl]{\strut{}{
    \begin{minipage}[h]{127.655896pt}
\textcolor{inkcol1}{\large{$1$}}\\
\end{minipage}}}}}%


\definecolor{inkcol1}{rgb}{0.0,0.0,0.0}
   \put(106.395467,120.847476){\rotatebox{360.0}{\makebox(0,0)[tl]{\strut{}{
    \begin{minipage}[h]{127.655896pt}
\textcolor{inkcol1}{\large{$2$}}\\
\end{minipage}}}}}%


\definecolor{inkcol1}{rgb}{0.0,0.0,0.0}
   \put(196.460937,45.072976){\rotatebox{360.0}{\makebox(0,0)[tl]{\strut{}{
    \begin{minipage}[h]{127.655896pt}
\textcolor{inkcol1}{\large{$k$}}\\
\end{minipage}}}}}%


\definecolor{inkcol1}{rgb}{0.0,0.0,0.0}
   \put(162.618107,102.997396){\rotatebox{360.0}{\makebox(0,0)[tl]{\strut{}{
    \begin{minipage}[h]{127.655896pt}
\textcolor{inkcol1}{\large{$2k\!-\!1$}}\\
\end{minipage}}}}}%


\definecolor{inkcol1}{rgb}{0.0,0.0,0.0}
   \put(116.851437,88.046496){\rotatebox{360.0}{\makebox(0,0)[tl]{\strut{}{
    \begin{minipage}[h]{127.655896pt}
\textcolor{inkcol1}{\large{$k\!-\!1$}}\\
\end{minipage}}}}}%


\definecolor{inkcol1}{rgb}{0.0,0.0,0.0}
   \put(186.474887,62.523116){\rotatebox{360.0}{\makebox(0,0)[tl]{\strut{}{
    \begin{minipage}[h]{127.655896pt}
\textcolor{inkcol1}{\large{$k\!+\!1$}}\\
\end{minipage}}}}}%


\definecolor{inkcol1}{rgb}{0.0,0.0,0.0}
   \put(206.908497,17.034216){\rotatebox{360.0}{\makebox(0,0)[tl]{\strut{}{
    \begin{minipage}[h]{127.655896pt}
\textcolor{inkcol1}{\Large{$x$}}\\
\end{minipage}}}}}%


\definecolor{inkcol1}{rgb}{0.0,0.0,0.0}
   \put(8.515647,15.962786){\rotatebox{360.0}{\makebox(0,0)[tl]{\strut{}{
    \begin{minipage}[h]{127.655896pt}
\textcolor{inkcol1}{\Large{$x$}}\\
\end{minipage}}}}}%


\definecolor{inkcol1}{rgb}{0.0,0.0,0.0}
   \put(107.265647,184.059786){\rotatebox{360.0}{\makebox(0,0)[tl]{\strut{}{
    \begin{minipage}[h]{127.655896pt}
\textcolor{inkcol1}{\Large{$x$}}\\
\end{minipage}}}}}%

 \end{picture}
\endgroup
}
\captionsetup{width=10cm}
\caption{Basic generators, $n=2k$.}\label{fig:OneCuspBasicEven}
\end{minipage}				
\end{figure}


\subsubsection{Multiple cusps, self edges}
Pick a face with a self edge, and extend to a maximal tree with $1$-cycle $G$ as in~Figure~\ref{fig:TreeWithCycle}. Let $T=G\setminus\varepsilon_1$, and let   
$B_n$ denote the matrix formed by the columns of $\alpha^*$ corresponding to the face points shown in Figure~\ref{fig:OneCycleBasic} together with the edge points on $T$.
We order the generators $x_i\otimes e_j$ of $C_0^n$ as follows
\begin{equation}
x_1\otimes e_1,\dots, x_1\otimes e_{n-1},\quad x_2\otimes e_{n-1},\dots ,x_2\otimes e_1
\end{equation}
with a similar scheme for the other vertices. The following is an immediate consequence of the definition of $\alpha^*$.



\begin{lemma}
The matrix $B_n$ is given by
\begin{equation}
B_n=\left(\begin{array}{c|c}
I_{n-1}+R_{n-1}& \\
S_{n-1}&I_T\otimes\Z^{n-1}\\\cline{1-1}
 0 &
\end{array}\right),
\end{equation}
where $I_T\otimes\Z^{n-1}$ is the matrix obtained from $I_T$ by replacing each non-zero entry by $I_{n-1}$.\qed
\end{lemma}
\begin{corollary}
The determinant of $B_n$ is $\pm n$.
\end{corollary}
\begin{proof}
This follows from
\begin{equation}
\det(B_n)=\pm\det(\begin{pmatrix}I_{n-1}+R_{n-1}&I_{n-1}\\S_{n-1}&I_{n-1}\end{pmatrix}=\pm n,
\end{equation}
where the second equality follows from Lemmas~\ref{lemma:BlockDeterminant} and \ref{lemma:DeterminantIdentity}.
\end{proof}
\begin{figure}[htpb]
\centering
\begin{minipage}[c]{0.48\textwidth}
\scalebox{0.8}{\begingroup
 \setlength{\unitlength}{0.8pt}
 \begin{picture}(419.827,187.53966)
 \put(0,0){\includegraphics{figures_gen/OneCycleBasic_included}}

\definecolor{inkcol1}{rgb}{0.0,0.0,0.0}
   \put(-0.205343,21.642686){\rotatebox{360.0}{\makebox(0,0)[tl]{\strut{}{
    \begin{minipage}[h]{127.655896pt}
\textcolor{inkcol1}{\Large{$x_1$}}\\
\end{minipage}}}}}%


\definecolor{inkcol1}{rgb}{0.0,0.0,0.0}
   \put(95.658267,193.225666){\rotatebox{360.0}{\makebox(0,0)[tl]{\strut{}{
    \begin{minipage}[h]{127.655896pt}
\textcolor{inkcol1}{\Large{$x_1$}}\\
\end{minipage}}}}}%


\definecolor{inkcol1}{rgb}{0.0,0.0,0.0}
   \put(191.849327,20.175936){\rotatebox{360.0}{\makebox(0,0)[tl]{\strut{}{
    \begin{minipage}[h]{127.655896pt}
\textcolor{inkcol1}{\Large{$x_2$}}\\
\end{minipage}}}}}%


\definecolor{inkcol1}{rgb}{0.0,0.0,0.0}
   \put(78.019267,168.399576){\rotatebox{360.0}{\makebox(0,0)[tl]{\strut{}{
    \begin{minipage}[h]{127.655896pt}
\textcolor{inkcol1}{\large{$1$}}\\
\end{minipage}}}}}%


\definecolor{inkcol1}{rgb}{0.0,0.0,0.0}
   \put(89.694227,141.095706){\rotatebox{360.0}{\makebox(0,0)[tl]{\strut{}{
    \begin{minipage}[h]{127.655896pt}
\textcolor{inkcol1}{\large{$2$}}\\
\end{minipage}}}}}%


\definecolor{inkcol1}{rgb}{0.0,0.0,0.0}
   \put(102.925917,64.756666){\rotatebox{360.0}{\makebox(0,0)[tl]{\strut{}{
    \begin{minipage}[h]{127.655896pt}
\textcolor{inkcol1}{\large{$n\!-\!2$}}\\
\end{minipage}}}}}%


\definecolor{inkcol1}{rgb}{0.0,0.0,0.0}
   \put(113.590717,47.850206){\rotatebox{360.0}{\makebox(0,0)[tl]{\strut{}{
    \begin{minipage}[h]{127.655896pt}
\textcolor{inkcol1}{\large{$n\!-\!1$}}\\
\end{minipage}}}}}%


\definecolor{inkcol1}{rgb}{0.0,0.0,0.0}
   \put(186.771677,46.744906){\rotatebox{360.0}{\makebox(0,0)[tl]{\strut{}{
    \begin{minipage}[h]{127.655896pt}
\textcolor{inkcol1}{\large{$n$}}\\
\end{minipage}}}}}%


\definecolor{inkcol1}{rgb}{0.0,0.0,0.0}
   \put(177.396207,65.247636){\rotatebox{360.0}{\makebox(0,0)[tl]{\strut{}{
    \begin{minipage}[h]{127.655896pt}
\textcolor{inkcol1}{\large{$n\!+\!1$}}\\
\end{minipage}}}}}%


\definecolor{inkcol1}{rgb}{0.0,0.0,0.0}
   \put(131.490297,141.366736){\rotatebox{360.0}{\makebox(0,0)[tl]{\strut{}{
    \begin{minipage}[h]{127.655896pt}
\textcolor{inkcol1}{\large{$2n\!-\!3$}}\\
\end{minipage}}}}}%


\definecolor{inkcol1}{rgb}{0.0,0.0,0.0}
   \put(119.579697,168.019386){\rotatebox{360.0}{\makebox(0,0)[tl]{\strut{}{
    \begin{minipage}[h]{127.655896pt}
\textcolor{inkcol1}{\large{$2n\!-\!2$}}\\
\end{minipage}}}}}%


\definecolor{inkcol1}{rgb}{0.0,0.0,0.0}
   \put(34.705867,112.924186){\rotatebox{360.0}{\makebox(0,0)[tl]{\strut{}{
    \begin{minipage}[h]{127.655896pt}
\textcolor{inkcol1}{\Large{$\varepsilon_1$}}\\
\end{minipage}}}}}%


\definecolor{inkcol1}{rgb}{0.0,0.0,0.0}
   \put(152.388637,113.934326){\rotatebox{360.0}{\makebox(0,0)[tl]{\strut{}{
    \begin{minipage}[h]{127.655896pt}
\textcolor{inkcol1}{\Large{$\varepsilon_2$}}\\
\end{minipage}}}}}%

 \end{picture}
\endgroup
}
\captionsetup{width=10cm}
\caption{Basic generators, tree with $1$-cycle.}\label{fig:OneCycleBasic}
\end{minipage}	
\begin{minipage}[c]{0.48\textwidth}
\scalebox{0.8}{\begingroup
 \setlength{\unitlength}{0.8pt}
 \begin{picture}(366.29214,190.0125)
 \put(0,0){\includegraphics{figures_gen/ThreeCycleBasic_included}}

\definecolor{inkcol1}{rgb}{0.0,0.0,0.0}
   \put(-0.205344,21.929366){\rotatebox{360.0}{\makebox(0,0)[tl]{\strut{}{
    \begin{minipage}[h]{127.655896pt}
\textcolor{inkcol1}{\Large{$x_2$}}\\
\end{minipage}}}}}%


\definecolor{inkcol1}{rgb}{0.0,0.0,0.0}
   \put(98.017767,195.698506){\rotatebox{360.0}{\makebox(0,0)[tl]{\strut{}{
    \begin{minipage}[h]{127.655896pt}
\textcolor{inkcol1}{\Large{$x_1$}}\\
\end{minipage}}}}}%


\definecolor{inkcol1}{rgb}{0.0,0.0,0.0}
   \put(194.689837,29.589416){\rotatebox{360.0}{\makebox(0,0)[tl]{\strut{}{
    \begin{minipage}[h]{127.655896pt}
\textcolor{inkcol1}{\Large{$x_3$}}\\
\end{minipage}}}}}%


\definecolor{inkcol1}{rgb}{0.0,0.0,0.0}
   \put(129.354927,160.917186){\rotatebox{360.0}{\makebox(0,0)[tl]{\strut{}{
    \begin{minipage}[h]{127.655896pt}
\textcolor{inkcol1}{\large{$n\!-\!1$}}\\
\end{minipage}}}}}%


\definecolor{inkcol1}{rgb}{0.0,0.0,0.0}
   \put(95.508487,124.573426){\rotatebox{360.0}{\makebox(0,0)[tl]{\strut{}{
    \begin{minipage}[h]{127.655896pt}
\textcolor{inkcol1}{\large{$n\!-\!2$}}\\
\end{minipage}}}}}%


\definecolor{inkcol1}{rgb}{0.0,0.0,0.0}
   \put(83.875607,80.353336){\rotatebox{360.0}{\makebox(0,0)[tl]{\strut{}{
    \begin{minipage}[h]{127.655896pt}
\textcolor{inkcol1}{\large{$2$}}\\
\end{minipage}}}}}%


\definecolor{inkcol1}{rgb}{0.0,0.0,0.0}
   \put(69.458617,60.335406){\rotatebox{360.0}{\makebox(0,0)[tl]{\strut{}{
    \begin{minipage}[h]{127.655896pt}
\textcolor{inkcol1}{\large{$1$}}\\
\end{minipage}}}}}%


\definecolor{inkcol1}{rgb}{0.0,0.0,0.0}
   \put(150.430547,17.243666){\rotatebox{360.0}{\makebox(0,0)[tl]{\strut{}{
    \begin{minipage}[h]{127.655896pt}
\textcolor{inkcol1}{\large{$2n\!-\!1$}}\\
\end{minipage}}}}}%


\definecolor{inkcol1}{rgb}{0.0,0.0,0.0}
   \put(126.090987,17.348266){\rotatebox{360.0}{\makebox(0,0)[tl]{\strut{}{
    \begin{minipage}[h]{127.655896pt}
\textcolor{inkcol1}{\large{$2n$}}\\
\end{minipage}}}}}%


\definecolor{inkcol1}{rgb}{0.0,0.0,0.0}
   \put(75.528997,17.348276){\rotatebox{360.0}{\makebox(0,0)[tl]{\strut{}{
    \begin{minipage}[h]{127.655896pt}
\textcolor{inkcol1}{\large{$3n\!-\!2$}}\\
\end{minipage}}}}}%


\definecolor{inkcol1}{rgb}{0.0,0.0,0.0}
   \put(22.713907,17.113666){\rotatebox{360.0}{\makebox(0,0)[tl]{\strut{}{
    \begin{minipage}[h]{127.655896pt}
\textcolor{inkcol1}{\large{$3n\!-\!3$}}\\
\end{minipage}}}}}%


\definecolor{inkcol1}{rgb}{0.0,0.0,0.0}
   \put(15.932547,64.220646){\rotatebox{360.0}{\makebox(0,0)[tl]{\strut{}{
    \begin{minipage}[h]{127.655896pt}
\textcolor{inkcol1}{\large{$n$}}\\
\end{minipage}}}}}%


\definecolor{inkcol1}{rgb}{0.0,0.0,0.0}
   \put(6.024096,85.938926){\rotatebox{360.0}{\makebox(0,0)[tl]{\strut{}{
    \begin{minipage}[h]{127.655896pt}
\textcolor{inkcol1}{\large{$n\!+\!1$}}\\
\end{minipage}}}}}%


\definecolor{inkcol1}{rgb}{0.0,0.0,0.0}
   \put(29.048917,143.770146){\rotatebox{360.0}{\makebox(0,0)[tl]{\strut{}{
    \begin{minipage}[h]{127.655896pt}
\textcolor{inkcol1}{\large{$2n\!-\!3$}}\\
\end{minipage}}}}}%


\definecolor{inkcol1}{rgb}{0.0,0.0,0.0}
   \put(153.676427,113.876286){\rotatebox{360.0}{\makebox(0,0)[tl]{\strut{}{
    \begin{minipage}[h]{127.655896pt}
\textcolor{inkcol1}{\Large{$\varepsilon_1$}}\\
\end{minipage}}}}}%


\definecolor{inkcol1}{rgb}{0.0,0.0,0.0}
   \put(35.488577,113.876286){\rotatebox{360.0}{\makebox(0,0)[tl]{\strut{}{
    \begin{minipage}[h]{127.655896pt}
\textcolor{inkcol1}{\Large{$\varepsilon_2$}}\\
\end{minipage}}}}}%


\definecolor{inkcol1}{rgb}{0.0,0.0,0.0}
   \put(96.855347,43.670686){\rotatebox{360.0}{\makebox(0,0)[tl]{\strut{}{
    \begin{minipage}[h]{127.655896pt}
\textcolor{inkcol1}{\Large{$\varepsilon_3$}}\\
\end{minipage}}}}}%


\definecolor{inkcol1}{rgb}{0.0,0.0,0.0}
   \put(41.746477,167.219866){\rotatebox{360.0}{\makebox(0,0)[tl]{\strut{}{
    \begin{minipage}[h]{127.655896pt}
\textcolor{inkcol1}{\large{$2n\!-\!2$}}\\
\end{minipage}}}}}%

 \end{picture}
\endgroup
}
\captionsetup{width=10cm}
\caption{Basic generators, tree with $3$-cycle.}\label{fig:ThreeCycleBasic}
\end{minipage}					
\end{figure}

\subsubsection{Multiple cusps, no self edge}
Pick a maximal tree with $3$-cycle $G$, and let $C_n$ be the matrix formed by the columns of $\alpha^*$ corresponding to the face points in Figure~\ref{fig:ThreeCycleBasic} together with the edge points on $T=G\setminus\varepsilon_1$.
\begin{lemma}
The matrix $C_n$ is given by
\begin{equation}
C_n=\left(\begin{array}{c|c}
I_{n-1}& \\
S_{n-1}&I_T\otimes\Z^{n-1}\\
R_{n-1}&\\\cline{1-1}
 0 &
\end{array}\right),
\end{equation}
\end{lemma}
\begin{corollary}
The determinant of $C_n$ is $\pm n$.
\end{corollary}
\begin{proof}
We have
\begin{equation}
\det(C_n)=\pm\det(M),\qquad M=\begin{pmatrix}I_{n-1}&I_{n-1}&\\S_{n-1}&I_{n-1}&I_{n-1}\\R_{n-1}&&I_{n-1}\end{pmatrix}.
\end{equation}
Using Lemma~\ref{lemma:BlockDeterminant} with $A=\left(\begin{smallmatrix}I_{n-1}&I_{n-1}\\S_{n-1}&I_{n-1}\end{smallmatrix}\right)$, $B=\left(\begin{smallmatrix}0\\I_{n-1}\end{smallmatrix}\right)$, $C=\left(\begin{smallmatrix}R_{n-1}&0\end{smallmatrix}\right)$, $D=I_{n-1}$, we have
\begin{equation}
\det(M)=\det\begin{pmatrix}I_{n-1}&I_{n-1}\\S_{n-1}-R_{n-1}&I_{n-1}\end{pmatrix}=\det(I_{n-1}+R_{n-1}-S_{n-1})=n,
\end{equation}
where the second equation follows from Lemma~\ref{lemma:BlockDeterminant}, and the third from Lemma~\ref{lemma:DeterminantIdentity}.
\end{proof}
This concludes the proof that $\alpha^*$ is basic.
%\subsection{The Ptolemy field}
%\label{sub.pfield}
%
%In this section we will define the Ptolemy field using some combinatorial
%results of gluing equations presented in \cite[Sec.4.6]{GHRS}.
%
%
%
%
%
%
%Fix an ideal triangulation $\CT$ of a compact oriented manifold $M$ with
%$t$ simplices and $r$ boundary torii components. Let $G$ denote the 
%1-skeleton of $\CT$ together with one ideal vertex for each cusp of $M$.
%$G$ is a finite graph with $r$ vertice and $t$ edges with incidence
%matrix $R$, an $r \times t$ matrix whose $(h,i)$ entry $r_{h,i}$
%gives the number of ends of the edge $i$ on the cusp $h$.
%
%$\CT$ has $t$ edges $e_i$ for $i=1,\dots,t$. For every edge $e_i$ of $\CT$ 
%consider the vector $E_i \in \BZ^{2t}$ of coefficients of the gluing equation 
%around $e_i$. Let $\Lambda \subset \BZ^{2t}$ denote the lattice of $\BZ$-linear 
%combinations of $E_i$ for $i=1,\dots,t$. From the work of Neumann-Zagier and 
%Neumann \cite{NeumannZagier,NeumannComb}, it follows that the matrix $R$ 
%gives the linear relations between the edge equations vectors $E_e$. More 
%presicely, we have
%$$
%\sum_i r_{hi} E_i=0 \quad \text{for all} \quad h=1,\dots,r \,.
%$$
%
%\begin{definition}
%\cite[Defn.4.2]{GHRS}
%A subset of the edges of a graph $\Gamma$ is a \emph{maximal tree with 1- 
%or 3-cycle} in $\Gamma$ if (together with the vertices) it consists of any 
%maximal tree $T$ together with one additional edge that either (1) is a loop 
%at one vertex, or (2) forms a 3-cycle together with two edges in $T$.
%\end{definition}
%
%\begin{theorem} 
%\cite[Thm.4.3]{GHRS}
%\label{thm.GHRS}
%There exists an integer basis for $\Lambda$ consisting of $t-r$ 
%of the edge equation coefficients $E_1, \ldots, E_n$.  In fact, we can choose
%such a basis by omitting $r$ edge equations corresponding to a maximal tree 
%with 1- or 3-cycle in $G$.
%\end{theorem}
%
%\begin{corollary}
%\label{cor.GHRS}
%Maximal trees $T$ with 1- or 3-cycle in $G$ always exist.
%If $T$ is such a tree with 1- or 3-cycles, let $v_1,\dots,v_{t-r}$ denote 
%the corresponding basis for $\Ker(R)$. 
%\end{corollary}
%
%Now if $x=(x_1,\dots,x_r) \in (\BC^*)^r$ and $c=(c_1,\dots,c_t)$ and
%$v \in \BZ^t$, let $c^v=\prod_{i=1}^t c_i^{v_i}$. It follows that
%$$
%x \cdot c^v = \prod_{i=1}^r x_i^{(R v)_i} c^v
%$$
%where $(R v)_i$ is the $i$th coodrinate of the vector $Rv \in \BZ^r$.
%Hence $c^v$ is invariant under $x$ if and only if $v$ is in the kernel of
%$R$. Corollary \ref{thm.GHRS} and the above discussion imply that
%\begin{itemize}
%\item
%the field 
%$$
%E_c=\BQ(c^{v_1},\dots,c^{v_{t-r}})
%$$
%is invariant under the action of $(\BC^*)^r$,
%\item
%independent of the choice of $T$.
%\end{itemize}
%This concludes the definition of the Ptolemy field.
%
%
%\begin{example} 
%\cite[Example 4.5]{GHRS}
%Let $M$ be the Whitehead link complement (with the triangulation given by 
%m129 in SnapPy notation). Then the matrix of edge equations is
%$$
%(A \mid B) = 
%\begin{pmatrix}
%2 & -1 & 1 & 1 & \vline & 1 & -2 &  0 & 0 \\
%-1 & 0 & 0 & 0  & \vline & -1 & 1 & 1 & 1 \\
%0 & 1 & -1 & -1 & \vline & 1 & 0 & -2 & -2 \\
%-1 & 0 & 0 & 0 & \vline & -1 & 1 & 1 & 1
%\end{pmatrix}
%$$
%The cusp incidence matrix is 
%$$
%R = \begin{pmatrix} 1 & 2 & 1 & 0 \\ 1 & 0 & 1 & 2 \end{pmatrix}
%$$
%and the corresponding graph $G$ is shown in Figure 
%\ref{wh_link_edges_and_cusps}.
%
%\begin{figure}[htbp]
%\begingroup
 \setlength{\unitlength}{0.8pt}
 \begin{picture}(232.4025,96.17975)
 \put(0,0){\includegraphics{figures_gen/wh_link_edges_and_cusps_included}}
 \end{picture}
\endgroup

%\caption{The graph of edges and cusps for the Whitehead link.}
%\label{wh_link_edges_and_cusps}
%\end{figure}
%
%The rows $E_0, E_1, E_2, E_3$ of $( A \mid  B)$ satisfy the relations: 
%$$
%E_0 + 2 E_1 + E_2   = 0 \text{ (from cusp 0)}
%$$
%and 
%$$
%E_0 + E_2  + 2 E_3  = 0 \text{ (from cusp 1)},
%$$
%which imply that $E_0=-2E_1-E_2$ and $E_1=E_3$. 
%The corresponding basis of $\Ker(R)$ is given by ${v_0,v_1}$ where
%$$
%\left(\begin{matrix}
%E_0 \\ E_1 \\ E_2 \\ E_3 \end{matrix}\right)=
%\left(\begin{matrix}
%-E_2-2E_3 \\ E_3 \\ E_2 \\ E_3 \end{matrix}\right)=
%E_2 v_1 + E_3 v_2, \qquad
%v_1=\left(\begin{matrix}
%-1 \\ 0 \\ 1 \\ 0 \end{matrix}\right), \qquad
%v_2=\left(\begin{matrix}
%-2 \\ 1 \\ 0 \\ 1 \end{matrix}\right)
%$$
%If we choose the tree $T$ with 1- or 3-cycles which consists of the edges 
%${e_0,e_1}$, then $E_2, E_3$ are linearly independent and form an integer 
%basis for the $\BZ$-span of $\{ E_0,E_1,E_ 2,E_3 \}$. 
%
%The action of $x=(x_0,x_1) \in (\BC^*)^2$ on $c=(c_0,c_1,c_2,c_3)$ is given
%by
%$$
%x \cdot c =(x_0 x_1 c_0, x_0^2 c_1, x_0 x_1 c_2, x_1^2 c_3)
%$$
%It follows that the following Ptolemy monomials
%$$
%(c^{v_0}, c^{v_1})=(c_0^{-1} c_2, c_0^{-2} c_1 c_3)
%$$
%are invariant under the $x$-action. The Ptolemy field is
%$E_c=\BQ(c_0^{-1} c_2, c_0^{-2} c_1 c_3)$.
%
%To show that $E_c$ is a subfield of the trace field $k_c$, consider the loop
%$(c_0c_2^{-1})$ {\bf figure missing}
%formed by starting on the vertex of $e_0$ on the $0$-cusp, 
%walk along $c_0$ to the $1$-cusp, join by a short path to the vertex of 
%$e_2$ on the $1$-cusp, then walk back on $e_2$ to the $0$-cusp and join 
%by a chort path on the $0$-cusp to the vertex of $e_0$. We have:
%%%%% see Mathematica file: zickert/LongShortEdges.Ptolemy.nb
%$$
%\tr( A(c_0) B(x_3) A(c_2)^{-1} B(x_1))=-c_0 c_2 x_1 x_3 \in k_c.
%$$
%Now we can alter the short paths of edges by adding peripheral curves. 
%The effect of this replaces $B(x_i)$ by $B(x_i+m_i)$ for $i=1,3$
%where we can assume that $m_i \neq 0$ (since the restriction of $\rho_c$ 
%to each boundary component is nontrivial).
%$$
%\Delta_{x,m}(f(x))=f(x+m,y)-f(x,y)
%$$ 
%It follows that
%$$
%\Delta_{x_3,m_3} \Delta_{x_1,m_1}(-c_0 c_2 x_1 x_3)= -c_0 c_2 m_1 m_3 \in k_c
%$$
%Likewise, the loop $(c_0c_0^{-1})$ gives that
%$$
%\tr( A(c_0) B(x_3) A(c_0)^{-1} B(x_1))=2-c_0^2 x_1 x_3 \in k_c
%$$
%and consequently that
%$$
%\Delta_{x_3,m_3} \Delta_{x_1,m_1}(2-c_0^2 x_1 x_3)=-c_0^2 m_1 m_3 \in k_c \,.
%$$
%It follows that 
%\begin{equation}
%\lbl{eq.c02}
%c_0^{-1} c_2, c_0^2 m_1 m_3 \in k_c \,.
%\end{equation}
%Now consider the loop $(c_1 c_0 c_3 c_0^{-1}$. Let
%It follows that
%\begin{multline}
%\tr( A(c_1) B(x_1) A(c_0) B(x_3) A(c_3) B(y_3) A(c_0)^{-1} B(y_1))= \\
%-c_0^2 c_1 c_3 x_1 x_3 y_1 y_3+\frac{c_0^2 c_1 x_1 y_1}{c_3}
%+\frac{c_0^2 c_3 x_3 y_3}{c_1}-\frac{c_0^2}{c_1 c_3}
%-\frac{c_1 c_3}{c_0y^2}+c_1 c_3 x_1 x_3+c_1 c_3 y_1 y_3
%\end{multline}
%It follows that
%\begin{multline}
%\Delta_{x_3,m_3} \Delta_{y_1,m_1}\Delta_{x_3,m_3} \Delta_{y_3,m_3}(
%\tr( A(c_1) B(x_1) A(c_0) B(x_3) A(c_3) B(y_3) A(c_0)^{-1} B(y_1))) = \\
%-c_0^2 c_1 c_3 m_1^2 m_3^2 \in k_c \,.
%\end{multline}
%Combined with \eqref{eq.c02}, it implies that
%$c_0^{-1} c_2, c_0^{-2} c_1 c_3 \in k_c$. This proves that $E_c$ is a subfield
%of $k_c$. The opposite conclusion holds since the representation $\rho_c$ is
%defined over $E_c$. So $E_c=k_c$.
%\end{example}
%
%\subsection{Proof of Theorem \ref{thm.1}}
%\label{sub.thm1}
%
%For simplicity, we consider the case of 1-cusp. Fix $c \in V(\CT)$ and
%consider the corresponding $\SL(2,\BC)$-representation $\rho_c$. Since
%$\rho_c$ is defined over $\BQ(c)$, it follows that $k_c$ is a subfield of
%$\BQ(c)$. Since $\rho_{x \cdot c}$ is conjugate to $\rho_c$, it follows that
%they have equal trace fields, hence $k_c$ is a subfield of the intersection
%$\cap_{x \in \BC^*} \BQ(x \cdot c)=E_c$.
%
%For the opposite inclusion, fix an edge $e$ of $\CT$ and choose a path
%$p$ of short edges that connect its two ends on the same boundary component
%of $M$. Consider the closed loop $e \cdot p$ on $\pi_1(M)$. Let $\rho_c(p)=
%B(x)$. Since the restriction of $\rho_c$ to $\pi_1(\pt M)$ is nontrivial, after
%a possible modification of $p$, we may assume that $x \neq 0$.
%Then, we have
%$$
%\tr(e \cdot p)=\tr(A(c)B(x))=c x \,.
%$$
%It follows that there exists $x \neq 0$ such that $c x \in k_c$ for 
%all Ptolemy coordinates $c$. Hence $c_e/c_{e'} \in k_c$ for all edges
%$e$ and $e'$ of $\CT$. The result follows. 
%
%It remains to deal with the multi-cusped case {\bf find out}.
%\qed
%
%%%%%%%%%%%%%%%%%%%%%%%%%%%%%%%%%%%%%%%%%%%%%%%%%%%%%%%%%%%%%%%%%%%%%%%%%%%%%%
%%%%%%%%%%%%%%%%%%%%%%%%%%%%%%%%%%%%%%%%%%%%%%%%%%%%%%%%%%%%%%%%%%%%%%%%%%%%%%



%%%%%%%%%%%%%%%%%%%%%%%%%%%%%%%%%%%%%%%%%%%%%%%%%%%%%%%%%%%%%%%%%%%%%%%%%%%%%
%%%%%%%%%%%%%%%%%%%%%%%%%%%%%%%%%%%%%%%%%%%%%%%%%%%%%%%%%%%%%%%%%%%%%%%%%%%%%

\bibliographystyle{plain}
\bibliography{BibFile}

\end{document}


